% Section 11 — Into the land: Jordan, Jericho, Ebal, Shiloh.
% Douay wording per build/ark-research/w2-rights.md §§1.8–1.10, 1.26, 1.29;
% corpus per ot-corpus.md §§19–25; archaeology per journey.md §§1.2–1.5.
\sectionguard
\section{Into the land: Jordan, Jericho, Ebal, Silo}

The chapters that carry the ark into the land are the ones in which it looks
most like the thing the rest of its story will deny it is. A river halts
because a box is carried into it. A walled city falls because a box is
carried round it. A reader who concludes that Israel possessed a weapon has
concluded nothing worse than what Israel's own elders will one day conclude,
with results this story will not spare him. The misreading is not a modern
impiety; it is the oldest mistake in the story. The book of Josue (the
Douay's name for Joshua) never rebukes the impression directly. It does
something quieter. It hedges its wonders about with distance, title, and
order, so that the careful reader comes away having watched not a weapon at
work but a King arriving.

% Drawing: the ark's journey, Sinai to Jerusalem. Sites are projected from
% their traditional identifications (x ~ longitude, y ~ latitude); the caption
% carries the numbered stages so the map itself stays a map.
\begin{sketchfig}
\begin{tikzpicture}[scale=0.88,
  march/.style={draw=black, line width=0.68pt, dash pattern=on 4.2pt off 2.5pt},
  stage/.style={circle, draw=black, line width=0.4pt, inner sep=0.75pt,
    font=\scriptsize, fill=white}]
% ---- neatline ---------------------------------------------------------------
\draw[pline] (0,-0.1) rectangle (11.2,12.05);
\begin{scope}
\clip (0,-0.1) rectangle (11.2,12.05);
% ---- the Great Sea: coast and waterlines ------------------------------------
\draw[pbold] (0.55,-0.1)
  .. controls (1.1,1.3) and (1.5,2.35) .. (2.22,3.72)
  .. controls (2.6,4.5) and (3.0,5.2) .. (3.30,6.00)
  .. controls (3.55,6.75) and (3.85,7.9) .. (4.14,8.70)
  .. controls (4.35,9.3) and (4.45,10.1) .. (4.62,10.68)
  .. controls (4.75,11.05) and (5.0,11.35) .. (5.5,12.05);
\draw[pwater, draw=black!40] (0.2,0.4)
  .. controls (0.8,1.6) and (1.25,2.6) .. (1.9,3.85)
  .. controls (2.3,4.6) and (2.7,5.3) .. (3.0,6.1)
  .. controls (3.25,6.85) and (3.55,7.95) .. (3.84,8.78)
  .. controls (4.05,9.38) and (4.15,10.15) .. (4.32,10.75)
  .. controls (4.45,11.1) and (4.7,11.45) .. (5.1,12.05);
\draw[pwater, draw=black!25] (0.0,1.15)
  .. controls (0.55,2.1) and (0.95,2.9) .. (1.55,4.0)
  .. controls (1.95,4.75) and (2.4,5.45) .. (2.7,6.25)
  .. controls (2.95,7.0) and (3.25,8.05) .. (3.54,8.88)
  .. controls (3.68,9.3) and (3.76,9.6) .. (3.8,9.85);
\draw[pwater, draw=black!15] (-0.35,1.9)
  .. controls (0.2,2.7) and (0.6,3.3) .. (1.2,4.3)
  .. controls (1.6,5.0) and (2.1,5.7) .. (2.4,6.45)
  .. controls (2.65,7.2) and (2.95,8.2) .. (3.24,9.0)
  .. controls (3.45,9.55) and (3.55,10.25) .. (3.72,10.85);
\node[stiny, rotate=68, black!60] at (1.35,6.2) {\scshape the great sea};
% ---- the Salt Sea -----------------------------------------------------------
\draw[pline] (7.80,4.32)
  .. controls (7.6,3.9) and (7.45,3.1) .. (7.35,2.4)
  .. controls (7.28,1.85) and (7.5,1.15) .. (7.85,0.72)
  .. controls (8.05,0.5) and (8.25,0.55) .. (8.3,0.85)
  .. controls (8.36,1.35) and (8.28,1.9) .. (8.22,2.35)
  .. controls (8.42,2.8) and (8.45,3.5) .. (8.32,3.95)
  .. controls (8.22,4.25) and (7.95,4.45) .. (7.80,4.32);
\begin{scope}
  \clip (7.80,4.32)
    .. controls (7.6,3.9) and (7.45,3.1) .. (7.35,2.4)
    .. controls (7.28,1.85) and (7.5,1.15) .. (7.85,0.72)
    .. controls (8.05,0.5) and (8.25,0.55) .. (8.3,0.85)
    .. controls (8.36,1.35) and (8.28,1.9) .. (8.22,2.35)
    .. controls (8.42,2.8) and (8.45,3.5) .. (8.32,3.95)
    .. controls (8.22,4.25) and (7.95,4.45) .. (7.80,4.32);
  \draw[draw=black!24, line width=0.16pt] (7.27,0.70) -- (7.69,0.70);
  \draw[draw=black!24, line width=0.16pt] (7.98,0.70) -- (8.20,0.70);
  \draw[draw=black!24, line width=0.16pt] (7.38,0.92) -- (7.74,0.92);
  \draw[draw=black!24, line width=0.16pt] (8.27,0.92) -- (8.50,0.92);
  \draw[draw=black!24, line width=0.16pt] (7.26,1.14) -- (7.52,1.14);
  \draw[draw=black!24, line width=0.16pt] (7.98,1.14) -- (8.20,1.14);
  \draw[draw=black!24, line width=0.16pt] (7.30,1.36) -- (7.53,1.36);
  \draw[draw=black!24, line width=0.16pt] (8.09,1.36) -- (8.37,1.36);
  \draw[draw=black!24, line width=0.16pt] (7.30,1.58) -- (7.57,1.58);
  \draw[draw=black!24, line width=0.16pt] (7.97,1.58) -- (8.24,1.58);
  \draw[draw=black!24, line width=0.16pt] (7.36,1.80) -- (7.55,1.80);
  \draw[draw=black!24, line width=0.16pt] (8.13,1.80) -- (8.31,1.80);
  \draw[draw=black!24, line width=0.16pt] (7.27,2.02) -- (7.55,2.02);
  \draw[draw=black!24, line width=0.16pt] (7.99,2.02) -- (8.33,2.02);
  \draw[draw=black!24, line width=0.16pt] (7.46,2.24) -- (7.68,2.24);
  \draw[draw=black!24, line width=0.16pt] (7.99,2.24) -- (8.20,2.24);
  \draw[draw=black!24, line width=0.16pt] (7.34,2.46) -- (7.56,2.46);
  \draw[draw=black!24, line width=0.16pt] (8.08,2.46) -- (8.48,2.46);
  \draw[draw=black!24, line width=0.16pt] (7.33,2.68) -- (7.58,2.68);
  \draw[draw=black!24, line width=0.16pt] (7.90,2.68) -- (8.30,2.68);
  \draw[draw=black!24, line width=0.16pt] (7.36,2.90) -- (7.61,2.90);
  \draw[draw=black!24, line width=0.16pt] (7.94,2.90) -- (8.15,2.90);
  \draw[draw=black!24, line width=0.16pt] (7.27,3.12) -- (7.46,3.12);
  \draw[draw=black!24, line width=0.16pt] (7.94,3.12) -- (8.18,3.12);
  \draw[draw=black!24, line width=0.16pt] (7.35,3.34) -- (7.68,3.34);
  \draw[draw=black!24, line width=0.16pt] (8.23,3.34) -- (8.50,3.34);
  \draw[draw=black!24, line width=0.16pt] (7.48,3.56) -- (7.72,3.56);
  \draw[draw=black!24, line width=0.16pt] (8.27,3.56) -- (8.46,3.56);
  \draw[draw=black!24, line width=0.16pt] (7.36,3.78) -- (7.77,3.78);
  \draw[draw=black!24, line width=0.16pt] (8.05,3.78) -- (8.25,3.78);
  \draw[draw=black!24, line width=0.16pt] (7.42,4.00) -- (7.68,4.00);
  \draw[draw=black!24, line width=0.16pt] (8.05,4.00) -- (8.25,4.00);
  \draw[draw=black!24, line width=0.16pt] (7.41,4.22) -- (7.71,4.22);
  \draw[draw=black!24, line width=0.16pt] (8.24,4.22) -- (8.50,4.22);
\end{scope}
\node[stiny, rotate=-78, black!70, fill=white, fill opacity=0.8, text opacity=1, inner sep=0.8pt] at (7.83,2.05) {the Salt Sea};
% ---- the Sea of Galilee -----------------------------------------------------
\draw[pline] (8.35,11.0)
  .. controls (8.65,10.95) and (8.78,10.6) .. (8.68,10.3)
  .. controls (8.6,10.05) and (8.45,9.9) .. (8.33,9.92)
  .. controls (8.12,9.98) and (8.05,10.4) .. (8.12,10.7)
  .. controls (8.17,10.92) and (8.22,11.02) .. (8.35,11.0);
\begin{scope}
  \clip (8.35,11.0)
    .. controls (8.65,10.95) and (8.78,10.6) .. (8.68,10.3)
    .. controls (8.6,10.05) and (8.45,9.9) .. (8.33,9.92)
    .. controls (8.12,9.98) and (8.05,10.4) .. (8.12,10.7)
    .. controls (8.17,10.92) and (8.22,11.02) .. (8.35,11.0);
  \draw[draw=black!24, line width=0.16pt] (8.25,9.95) -- (8.55,9.95);
  \draw[draw=black!24, line width=0.16pt] (8.09,10.15) -- (8.31,10.15);
  \draw[draw=black!24, line width=0.16pt] (8.62,10.15) -- (8.85,10.15);
  \draw[draw=black!24, line width=0.16pt] (8.01,10.35) -- (8.33,10.35);
  \draw[draw=black!24, line width=0.16pt] (8.73,10.35) -- (8.85,10.35);
  \draw[draw=black!24, line width=0.16pt] (8.23,10.55) -- (8.61,10.55);
  \draw[draw=black!24, line width=0.16pt] (8.05,10.75) -- (8.34,10.75);
  \draw[draw=black!24, line width=0.16pt] (8.70,10.75) -- (8.85,10.75);
  \draw[draw=black!24, line width=0.16pt] (8.09,10.95) -- (8.38,10.95);
\end{scope}
\node[stiny, anchor=west, black!70] at (8.9,10.5) {Sea of\\Galilee};
% ---- the Jordan, meandering -------------------------------------------------
\draw[draw=black!75, line width=0.38pt] (8.3,9.92)
  .. controls (8.1,9.5) and (8.35,9.1) .. (8.15,8.7)
  .. controls (7.95,8.3) and (8.25,7.9) .. (8.05,7.5)
  .. controls (7.85,7.1) and (8.15,6.7) .. (7.95,6.3)
  .. controls (7.78,5.95) and (8.0,5.5) .. (7.88,5.1)
  .. controls (7.82,4.85) and (7.8,4.6) .. (7.80,4.35);
\node[stiny, rotate=-83, black!60] at (7.72,7.7) {Jordan};
% ---- terrain: hachured ridges ----------------------------------------------
\draw[draw=black!50, line width=0.17pt] (5.521,2.096) -- (5.375,2.155);
\draw[draw=black!56, line width=0.17pt] (5.555,2.083) -- (5.693,2.028);
\draw[draw=black!56, line width=0.17pt] (5.557,2.149) -- (5.462,2.187);
\draw[draw=black!57, line width=0.17pt] (5.579,2.126) -- (5.820,2.029);
\draw[draw=black!67, line width=0.17pt] (5.568,2.172) -- (5.459,2.216);
\draw[draw=black!61, line width=0.17pt] (5.604,2.156) -- (5.815,2.072);
\draw[draw=black!52, line width=0.17pt] (5.574,2.204) -- (5.349,2.294);
\draw[draw=black!63, line width=0.17pt] (5.618,2.187) -- (5.826,2.103);
\draw[draw=black!52, line width=0.17pt] (5.579,2.227) -- (5.361,2.315);
\draw[draw=black!51, line width=0.17pt] (5.614,2.211) -- (5.727,2.166);
\draw[draw=black!57, line width=0.17pt] (5.613,2.282) -- (5.395,2.369);
\draw[draw=black!63, line width=0.17pt] (5.626,2.261) -- (5.792,2.195);
\draw[draw=black!62, line width=0.17pt] (5.617,2.314) -- (5.435,2.386);
\draw[draw=black!65, line width=0.17pt] (5.657,2.314) -- (5.793,2.260);
\draw[draw=black!57, line width=0.17pt] (5.615,2.353) -- (5.442,2.422);
\draw[draw=black!61, line width=0.17pt] (5.647,2.329) -- (5.764,2.283);
\draw[draw=black!58, line width=0.17pt] (5.646,2.356) -- (5.480,2.423);
\draw[draw=black!60, line width=0.17pt] (5.687,2.369) -- (5.882,2.291);
\draw[draw=black!68, line width=0.17pt] (5.653,2.415) -- (5.484,2.482);
\draw[draw=black!59, line width=0.17pt] (5.705,2.408) -- (5.833,2.357);
\draw[draw=black!68, line width=0.17pt] (5.673,2.434) -- (5.448,2.523);
\draw[draw=black!62, line width=0.17pt] (5.700,2.441) -- (5.883,2.368);
\draw[draw=black!52, line width=0.17pt] (5.686,2.467) -- (5.478,2.551);
\draw[draw=black!55, line width=0.17pt] (5.709,2.468) -- (5.813,2.427);
\draw[draw=black!65, line width=0.17pt] (5.697,2.514) -- (5.461,2.609);
\draw[draw=black!67, line width=0.17pt] (5.736,2.478) -- (5.933,2.400);
\draw[draw=black!61, line width=0.17pt] (5.712,2.535) -- (5.604,2.579);
\draw[draw=black!57, line width=0.17pt] (5.737,2.546) -- (5.930,2.468);
\draw[draw=black!58, line width=0.17pt] (5.710,2.576) -- (5.560,2.636);
\draw[draw=black!55, line width=0.17pt] (5.748,2.576) -- (5.944,2.498);
\draw[draw=black!53, line width=0.17pt] (5.746,2.615) -- (5.554,2.692);
\draw[draw=black!50, line width=0.17pt] (5.779,2.603) -- (5.906,2.552);
\draw[draw=black!60, line width=0.17pt] (5.735,2.650) -- (5.542,2.727);
\draw[draw=black!68, line width=0.17pt] (5.775,2.620) -- (5.922,2.561);
\draw[draw=black!64, line width=0.17pt] (5.752,2.663) -- (5.514,2.758);
\draw[draw=black!62, line width=0.17pt] (5.788,2.649) -- (5.989,2.569);
\draw[draw=black!57, line width=0.17pt] (5.761,2.728) -- (5.574,2.803);
\draw[draw=black!54, line width=0.17pt] (5.816,2.714) -- (5.929,2.669);
\draw[draw=black!57, line width=0.17pt] (5.770,2.755) -- (5.568,2.835);
\draw[draw=black!56, line width=0.17pt] (5.834,2.732) -- (5.953,2.684);
\draw[draw=black!55, line width=0.17pt] (5.803,2.756) -- (5.612,2.833);
\draw[draw=black!66, line width=0.17pt] (5.825,2.752) -- (6.016,2.676);
\draw[draw=black!59, line width=0.17pt] (5.799,2.809) -- (5.613,2.883);
\draw[draw=black!68, line width=0.17pt] (5.844,2.784) -- (6.046,2.703);
\draw[draw=black!66, line width=0.17pt] (5.832,2.855) -- (5.631,2.935);
\draw[draw=black!50, line width=0.17pt] (5.857,2.828) -- (6.068,2.743);
\draw[draw=black!64, line width=0.17pt] (5.850,2.875) -- (5.668,2.947);
\draw[draw=black!59, line width=0.17pt] (5.873,2.851) -- (6.001,2.800);
\draw[draw=black!60, line width=0.17pt] (5.837,2.895) -- (5.661,2.965);
\draw[draw=black!56, line width=0.17pt] (5.876,2.866) -- (6.087,2.781);
\draw[draw=black!63, line width=0.17pt] (5.867,2.943) -- (5.701,3.010);
\draw[draw=black!61, line width=0.17pt] (5.890,2.921) -- (6.067,2.850);
\draw[draw=black!64, line width=0.17pt] (5.879,2.962) -- (5.765,3.007);
\draw[draw=black!51, line width=0.17pt] (5.917,2.961) -- (6.127,2.877);
\draw[draw=black!56, line width=0.17pt] (5.902,3.012) -- (5.735,3.079);
\draw[draw=black!50, line width=0.17pt] (5.937,2.998) -- (6.066,2.946);
\draw[draw=black!63, line width=0.17pt] (5.916,3.050) -- (5.749,3.117);
\draw[draw=black!58, line width=0.17pt] (5.936,3.014) -- (6.052,2.967);
\draw[draw=black!65, line width=0.17pt] (5.924,3.051) -- (5.722,3.132);
\draw[draw=black!56, line width=0.17pt] (5.940,3.054) -- (6.139,2.975);
\draw[draw=black!53, line width=0.17pt] (5.917,3.110) -- (5.745,3.179);
\draw[draw=black!51, line width=0.17pt] (5.974,3.088) -- (6.074,3.048);
\draw[draw=black!61, line width=0.17pt] (5.938,3.119) -- (5.790,3.151);
\draw[draw=black!50, line width=0.17pt] (5.977,3.135) -- (6.102,3.109);
\draw[draw=black!67, line width=0.17pt] (5.956,3.167) -- (5.788,3.202);
\draw[draw=black!56, line width=0.17pt] (5.971,3.166) -- (6.077,3.143);
\draw[draw=black!65, line width=0.17pt] (5.952,3.209) -- (5.711,3.259);
\draw[draw=black!61, line width=0.17pt] (5.984,3.199) -- (6.178,3.158);
\draw[draw=black!63, line width=0.17pt] (5.971,3.238) -- (5.834,3.267);
\draw[draw=black!52, line width=0.17pt] (5.999,3.223) -- (6.156,3.190);
\draw[draw=black!54, line width=0.17pt] (5.960,3.270) -- (5.784,3.307);
\draw[draw=black!63, line width=0.17pt] (6.000,3.277) -- (6.113,3.253);
\draw[draw=black!67, line width=0.17pt] (5.963,3.315) -- (5.801,3.350);
\draw[draw=black!55, line width=0.17pt] (6.023,3.302) -- (6.267,3.250);
\draw[draw=black!60, line width=0.17pt] (5.966,3.326) -- (5.729,3.376);
\draw[draw=black!53, line width=0.17pt] (6.030,3.324) -- (6.205,3.287);
\draw[draw=black!53, line width=0.17pt] (5.998,3.397) -- (5.759,3.447);
\draw[draw=black!53, line width=0.17pt] (6.017,3.371) -- (6.255,3.321);
\draw[draw=black!53, line width=0.17pt] (5.996,3.413) -- (5.791,3.456);
\draw[draw=black!50, line width=0.17pt] (6.020,3.395) -- (6.235,3.349);
\draw[draw=black!56, line width=0.17pt] (5.996,3.438) -- (5.829,3.474);
\draw[draw=black!61, line width=0.17pt] (6.052,3.460) -- (6.298,3.408);
\draw[draw=black!51, line width=0.17pt] (6.005,3.507) -- (5.780,3.554);
\draw[draw=black!61, line width=0.17pt] (6.063,3.494) -- (6.171,3.471);
\draw[draw=black!52, line width=0.17pt] (6.014,3.515) -- (5.831,3.554);
\draw[draw=black!67, line width=0.17pt] (6.049,3.516) -- (6.200,3.484);
\draw[draw=black!57, line width=0.17pt] (6.033,3.556) -- (5.855,3.593);
\draw[draw=black!68, line width=0.17pt] (6.054,3.572) -- (6.194,3.542);
\draw[draw=black!56, line width=0.17pt] (6.043,3.608) -- (5.864,3.645);
\draw[draw=black!63, line width=0.17pt] (6.070,3.579) -- (6.234,3.545);
\draw[draw=black!51, line width=0.17pt] (6.050,3.618) -- (5.849,3.660);
\draw[draw=black!63, line width=0.17pt] (6.086,3.607) -- (6.283,3.565);
\draw[draw=black!55, line width=0.17pt] (6.047,3.675) -- (5.892,3.708);
\draw[draw=black!68, line width=0.17pt] (6.098,3.666) -- (6.284,3.627);
\draw[draw=black!63, line width=0.17pt] (6.068,3.703) -- (5.912,3.736);
\draw[draw=black!50, line width=0.17pt] (6.108,3.702) -- (6.302,3.662);
\draw[draw=black!62, line width=0.17pt] (6.063,3.748) -- (5.898,3.782);
\draw[draw=black!58, line width=0.17pt] (6.117,3.721) -- (6.252,3.692);
\draw[draw=black!63, line width=0.17pt] (6.072,3.771) -- (5.964,3.793);
\draw[draw=black!52, line width=0.17pt] (6.098,3.770) -- (6.328,3.722);
\draw[draw=black!60, line width=0.17pt] (6.084,3.795) -- (5.939,3.825);
\draw[draw=black!65, line width=0.17pt] (6.112,3.796) -- (6.226,3.772);
\draw[draw=black!62, line width=0.17pt] (6.079,3.840) -- (5.827,3.893);
\draw[draw=black!53, line width=0.17pt] (6.134,3.820) -- (6.331,3.779);
\draw[draw=black!54, line width=0.17pt] (6.097,3.876) -- (5.966,3.904);
\draw[draw=black!55, line width=0.17pt] (6.121,3.876) -- (6.298,3.839);
\draw[draw=black!52, line width=0.17pt] (6.094,3.906) -- (5.863,3.955);
\draw[draw=black!64, line width=0.17pt] (6.139,3.904) -- (6.380,3.853);
\draw[draw=black!65, line width=0.17pt] (6.104,3.952) -- (5.849,4.005);
\draw[draw=black!63, line width=0.17pt] (6.143,3.946) -- (6.366,3.899);
\draw[draw=black!60, line width=0.17pt] (6.106,3.973) -- (5.914,4.014);
\draw[draw=black!59, line width=0.17pt] (6.138,3.951) -- (6.381,3.900);
\draw[draw=black!66, line width=0.17pt] (6.135,4.021) -- (6.018,4.045);
\draw[draw=black!58, line width=0.17pt] (6.153,4.008) -- (6.280,3.981);
\draw[draw=black!57, line width=0.17pt] (6.129,4.045) -- (6.015,4.069);
\draw[draw=black!68, line width=0.17pt] (6.179,4.046) -- (6.281,4.024);
\draw[draw=black!63, line width=0.17pt] (6.122,4.077) -- (5.958,4.058);
\draw[draw=black!67, line width=0.17pt] (6.176,4.091) -- (6.376,4.115);
\draw[draw=black!52, line width=0.17pt] (6.110,4.105) -- (5.976,4.089);
\draw[draw=black!68, line width=0.17pt] (6.167,4.110) -- (6.377,4.134);
\draw[draw=black!68, line width=0.17pt] (6.111,4.149) -- (5.926,4.128);
\draw[draw=black!65, line width=0.17pt] (6.168,4.152) -- (6.303,4.167);
\draw[draw=black!58, line width=0.17pt] (6.123,4.180) -- (5.992,4.164);
\draw[draw=black!58, line width=0.17pt] (6.162,4.170) -- (6.417,4.200);
\draw[draw=black!55, line width=0.17pt] (6.111,4.196) -- (5.976,4.180);
\draw[draw=black!65, line width=0.17pt] (6.150,4.213) -- (6.407,4.243);
\draw[draw=black!54, line width=0.17pt] (6.094,4.248) -- (5.981,4.234);
\draw[draw=black!56, line width=0.17pt] (6.153,4.261) -- (6.263,4.274);
\draw[draw=black!65, line width=0.17pt] (6.091,4.313) -- (5.972,4.299);
\draw[draw=black!51, line width=0.17pt] (6.150,4.296) -- (6.345,4.319);
\draw[draw=black!62, line width=0.17pt] (6.089,4.345) -- (5.839,4.315);
\draw[draw=black!50, line width=0.17pt] (6.130,4.325) -- (6.245,4.339);
\draw[draw=black!65, line width=0.17pt] (6.090,4.372) -- (5.943,4.355);
\draw[draw=black!52, line width=0.17pt] (6.147,4.377) -- (6.335,4.399);
\draw[draw=black!53, line width=0.17pt] (6.085,4.395) -- (5.962,4.381);
\draw[draw=black!68, line width=0.17pt] (6.135,4.387) -- (6.272,4.403);
\draw[draw=black!65, line width=0.17pt] (6.080,4.426) -- (5.915,4.406);
\draw[draw=black!67, line width=0.17pt] (6.111,4.441) -- (6.266,4.460);
\draw[draw=black!53, line width=0.17pt] (6.092,4.482) -- (5.945,4.465);
\draw[draw=black!65, line width=0.17pt] (6.120,4.462) -- (6.302,4.484);
\draw[draw=black!50, line width=0.17pt] (6.065,4.483) -- (5.837,4.456);
\draw[draw=black!50, line width=0.17pt] (6.103,4.513) -- (6.243,4.529);
\draw[draw=black!53, line width=0.17pt] (6.086,4.542) -- (5.912,4.522);
\draw[draw=black!52, line width=0.17pt] (6.099,4.545) -- (6.269,4.565);
\draw[draw=black!53, line width=0.17pt] (6.076,4.567) -- (5.874,4.543);
\draw[draw=black!60, line width=0.17pt] (6.100,4.578) -- (6.358,4.608);
\draw[draw=black!59, line width=0.17pt] (6.060,4.616) -- (5.898,4.597);
\draw[draw=black!54, line width=0.17pt] (6.097,4.611) -- (6.200,4.623);
\draw[draw=black!55, line width=0.17pt] (6.055,4.631) -- (5.919,4.615);
\draw[draw=black!62, line width=0.17pt] (6.097,4.638) -- (6.282,4.660);
\draw[draw=black!63, line width=0.17pt] (6.062,4.683) -- (5.960,4.671);
\draw[draw=black!57, line width=0.17pt] (6.109,4.665) -- (6.285,4.685);
\draw[draw=black!60, line width=0.17pt] (6.055,4.697) -- (5.806,4.668);
\draw[draw=black!56, line width=0.17pt] (6.085,4.703) -- (6.197,4.716);
\draw[draw=black!61, line width=0.17pt] (6.045,4.733) -- (5.941,4.721);
\draw[draw=black!63, line width=0.17pt] (6.084,4.739) -- (6.303,4.765);
\draw[draw=black!54, line width=0.17pt] (6.037,4.805) -- (5.836,4.781);
\draw[draw=black!65, line width=0.17pt] (6.075,4.791) -- (6.249,4.811);
\draw[draw=black!51, line width=0.17pt] (6.027,4.835) -- (5.834,4.812);
\draw[draw=black!54, line width=0.17pt] (6.086,4.823) -- (6.339,4.853);
\draw[draw=black!61, line width=0.17pt] (6.040,4.870) -- (5.879,4.851);
\draw[draw=black!54, line width=0.17pt] (6.078,4.863) -- (6.249,4.884);
\draw[draw=black!67, line width=0.17pt] (6.040,4.882) -- (5.939,4.870);
\draw[draw=black!60, line width=0.17pt] (6.080,4.899) -- (6.213,4.915);
\draw[draw=black!59, line width=0.17pt] (6.026,4.928) -- (5.869,4.962);
\draw[draw=black!68, line width=0.17pt] (6.087,4.909) -- (6.197,4.885);
\draw[draw=black!53, line width=0.17pt] (6.032,4.965) -- (5.809,5.015);
\draw[draw=black!55, line width=0.17pt] (6.077,4.951) -- (6.205,4.922);
\draw[draw=black!63, line width=0.17pt] (6.054,5.004) -- (5.823,5.055);
\draw[draw=black!53, line width=0.17pt] (6.076,4.983) -- (6.329,4.927);
\draw[draw=black!51, line width=0.17pt] (6.058,5.037) -- (5.855,5.083);
\draw[draw=black!53, line width=0.17pt] (6.100,5.041) -- (6.198,5.020);
\draw[draw=black!55, line width=0.17pt] (6.075,5.093) -- (5.938,5.123);
\draw[draw=black!61, line width=0.17pt] (6.109,5.063) -- (6.329,5.014);
\draw[draw=black!52, line width=0.17pt] (6.083,5.115) -- (5.891,5.158);
\draw[draw=black!65, line width=0.17pt] (6.117,5.111) -- (6.310,5.068);
\draw[draw=black!66, line width=0.17pt] (6.094,5.144) -- (5.892,5.189);
\draw[draw=black!55, line width=0.17pt] (6.126,5.143) -- (6.238,5.118);
\draw[draw=black!64, line width=0.17pt] (6.096,5.185) -- (5.891,5.231);
\draw[draw=black!51, line width=0.17pt] (6.128,5.184) -- (6.231,5.161);
\draw[draw=black!54, line width=0.17pt] (6.109,5.215) -- (5.939,5.253);
\draw[draw=black!63, line width=0.17pt] (6.129,5.197) -- (6.367,5.144);
\draw[draw=black!60, line width=0.17pt] (6.091,5.253) -- (5.970,5.280);
\draw[draw=black!68, line width=0.17pt] (6.131,5.225) -- (6.267,5.195);
\draw[draw=black!64, line width=0.17pt] (6.108,5.280) -- (5.999,5.305);
\draw[draw=black!67, line width=0.17pt] (6.152,5.286) -- (6.311,5.251);
\draw[draw=black!58, line width=0.17pt] (6.111,5.320) -- (5.873,5.373);
\draw[draw=black!52, line width=0.17pt] (6.170,5.298) -- (6.288,5.272);
\draw[draw=black!67, line width=0.17pt] (6.142,5.371) -- (5.892,5.427);
\draw[draw=black!63, line width=0.17pt] (6.174,5.349) -- (6.290,5.323);
\draw[draw=black!65, line width=0.17pt] (6.137,5.390) -- (6.037,5.412);
\draw[draw=black!56, line width=0.17pt] (6.186,5.387) -- (6.292,5.364);
\draw[draw=black!64, line width=0.17pt] (6.137,5.434) -- (5.989,5.467);
\draw[draw=black!57, line width=0.17pt] (6.196,5.427) -- (6.329,5.398);
\draw[draw=black!52, line width=0.17pt] (6.156,5.462) -- (5.999,5.497);
\draw[draw=black!61, line width=0.17pt] (6.199,5.453) -- (6.339,5.422);
\draw[draw=black!68, line width=0.17pt] (6.153,5.488) -- (6.008,5.521);
\draw[draw=black!53, line width=0.17pt] (6.201,5.488) -- (6.372,5.450);
\draw[draw=black!55, line width=0.17pt] (6.170,5.534) -- (6.060,5.558);
\draw[draw=black!60, line width=0.17pt] (6.217,5.530) -- (6.388,5.492);
\draw[draw=black!53, line width=0.17pt] (6.173,5.568) -- (6.056,5.594);
\draw[draw=black!58, line width=0.17pt] (6.214,5.554) -- (6.400,5.513);
\draw[draw=black!51, line width=0.17pt] (6.171,5.592) -- (5.954,5.640);
\draw[draw=black!68, line width=0.17pt] (6.212,5.576) -- (6.457,5.521);
\draw[draw=black!61, line width=0.17pt] (6.195,5.640) -- (6.023,5.678);
\draw[draw=black!54, line width=0.17pt] (6.232,5.628) -- (6.341,5.604);
\draw[draw=black!65, line width=0.17pt] (6.186,5.676) -- (5.937,5.731);
\draw[draw=black!51, line width=0.17pt] (6.251,5.653) -- (6.458,5.607);
\draw[draw=black!50, line width=0.17pt] (6.206,5.703) -- (6.013,5.746);
\draw[draw=black!52, line width=0.17pt] (6.248,5.692) -- (6.383,5.662);
\draw[draw=black!50, line width=0.17pt] (6.229,5.728) -- (6.092,5.759);
\draw[draw=black!58, line width=0.17pt] (6.245,5.726) -- (6.433,5.684);
\draw[draw=black!56, line width=0.17pt] (6.214,5.797) -- (6.077,5.827);
\draw[draw=black!64, line width=0.17pt] (6.257,5.763) -- (6.438,5.723);
\draw[draw=black!66, line width=0.17pt] (6.238,5.805) -- (6.094,5.841);
\draw[draw=black!60, line width=0.17pt] (6.275,5.802) -- (6.377,5.776);
\draw[draw=black!68, line width=0.17pt] (6.238,5.827) -- (5.986,5.890);
\draw[draw=black!58, line width=0.17pt] (6.272,5.824) -- (6.477,5.773);
\draw[draw=black!55, line width=0.17pt] (6.262,5.881) -- (6.043,5.936);
\draw[draw=black!54, line width=0.17pt] (6.300,5.870) -- (6.417,5.841);
\draw[draw=black!54, line width=0.17pt] (6.267,5.907) -- (6.135,5.940);
\draw[draw=black!59, line width=0.17pt] (6.306,5.885) -- (6.450,5.849);
\draw[draw=black!50, line width=0.17pt] (6.257,5.942) -- (6.045,5.995);
\draw[draw=black!56, line width=0.17pt] (6.309,5.920) -- (6.477,5.878);
\draw[draw=black!51, line width=0.17pt] (6.268,5.968) -- (6.050,6.023);
\draw[draw=black!59, line width=0.17pt] (6.324,5.955) -- (6.447,5.924);
\draw[draw=black!57, line width=0.17pt] (6.292,6.027) -- (6.109,6.073);
\draw[draw=black!54, line width=0.17pt] (6.337,6.027) -- (6.568,5.969);
\draw[draw=black!59, line width=0.17pt] (6.288,6.053) -- (6.147,6.088);
\draw[draw=black!66, line width=0.17pt] (6.320,6.057) -- (6.549,5.999);
\draw[draw=black!65, line width=0.17pt] (6.303,6.096) -- (6.200,6.121);
\draw[draw=black!50, line width=0.17pt] (6.328,6.077) -- (6.469,6.042);
\draw[draw=black!54, line width=0.17pt] (6.316,6.122) -- (6.140,6.166);
\draw[draw=black!57, line width=0.17pt] (6.349,6.113) -- (6.488,6.078);
\draw[draw=black!65, line width=0.17pt] (6.315,6.146) -- (6.206,6.174);
\draw[draw=black!67, line width=0.17pt] (6.354,6.158) -- (6.538,6.112);
\draw[draw=black!62, line width=0.17pt] (6.336,6.183) -- (6.120,6.237);
\draw[draw=black!61, line width=0.17pt] (6.363,6.166) -- (6.503,6.131);
\draw[draw=black!62, line width=0.17pt] (6.343,6.231) -- (6.147,6.280);
\draw[draw=black!57, line width=0.17pt] (6.363,6.213) -- (6.568,6.162);
\draw[draw=black!67, line width=0.17pt] (6.353,6.258) -- (6.123,6.316);
\draw[draw=black!53, line width=0.17pt] (6.393,6.252) -- (6.622,6.194);
\draw[draw=black!66, line width=0.17pt] (6.361,6.299) -- (6.143,6.353);
\draw[draw=black!56, line width=0.17pt] (6.380,6.281) -- (6.541,6.240);
\draw[draw=black!52, line width=0.17pt] (6.379,6.328) -- (6.185,6.377);
\draw[draw=black!58, line width=0.17pt] (6.391,6.335) -- (6.616,6.279);
\draw[draw=black!59, line width=0.17pt] (6.364,6.366) -- (6.252,6.394);
\draw[draw=black!62, line width=0.17pt] (6.423,6.351) -- (6.669,6.289);
\draw[draw=black!66, line width=0.17pt] (6.374,6.393) -- (6.146,6.451);
\draw[draw=black!54, line width=0.17pt] (6.430,6.389) -- (6.638,6.337);
\draw[draw=black!55, line width=0.17pt] (6.387,6.448) -- (6.221,6.490);
\draw[draw=black!54, line width=0.17pt] (6.426,6.443) -- (6.546,6.413);
\draw[draw=black!52, line width=0.17pt] (6.394,6.485) -- (6.190,6.536);
\draw[draw=black!59, line width=0.17pt] (6.445,6.476) -- (6.673,6.419);
\draw[draw=black!65, line width=0.17pt] (6.420,6.507) -- (6.277,6.543);
\draw[draw=black!66, line width=0.17pt] (6.450,6.500) -- (6.679,6.443);
\draw[draw=black!52, line width=0.17pt] (6.420,6.548) -- (6.187,6.606);
\draw[draw=black!56, line width=0.17pt] (6.466,6.530) -- (6.608,6.494);
\draw[draw=black!50, line width=0.17pt] (6.409,6.582) -- (6.273,6.616);
\draw[draw=black!64, line width=0.17pt] (6.473,6.556) -- (6.718,6.495);
\draw[draw=black!58, line width=0.17pt] (6.427,6.611) -- (6.227,6.640);
\draw[draw=black!66, line width=0.17pt] (6.484,6.596) -- (6.692,6.566);
\draw[draw=black!67, line width=0.17pt] (6.445,6.637) -- (6.334,6.653);
\draw[draw=black!59, line width=0.17pt] (6.476,6.637) -- (6.665,6.610);
\draw[draw=black!50, line width=0.17pt] (6.440,6.665) -- (6.185,6.702);
\draw[draw=black!62, line width=0.17pt] (6.478,6.662) -- (6.624,6.641);
\draw[draw=black!59, line width=0.17pt] (6.455,6.690) -- (6.208,6.725);
\draw[draw=black!68, line width=0.17pt] (6.495,6.691) -- (6.752,6.654);
\draw[draw=black!50, line width=0.17pt] (6.440,6.744) -- (6.214,6.776);
\draw[draw=black!68, line width=0.17pt] (6.480,6.745) -- (6.724,6.710);
\draw[draw=black!50, line width=0.17pt] (6.454,6.775) -- (6.230,6.807);
\draw[draw=black!60, line width=0.17pt] (6.487,6.775) -- (6.665,6.749);
\draw[draw=black!61, line width=0.17pt] (6.457,6.805) -- (6.246,6.835);
\draw[draw=black!60, line width=0.17pt] (6.492,6.808) -- (6.672,6.782);
\draw[draw=black!58, line width=0.17pt] (6.479,6.863) -- (6.326,6.885);
\draw[draw=black!59, line width=0.17pt] (6.493,6.858) -- (6.699,6.828);
\draw[draw=black!51, line width=0.17pt] (6.484,6.874) -- (6.271,6.905);
\draw[draw=black!51, line width=0.17pt] (6.525,6.871) -- (6.628,6.856);
\draw[draw=black!51, line width=0.17pt] (6.488,6.923) -- (6.366,6.941);
\draw[draw=black!55, line width=0.17pt] (6.506,6.919) -- (6.724,6.888);
\draw[draw=black!56, line width=0.17pt] (6.491,6.978) -- (6.316,7.003);
\draw[draw=black!66, line width=0.17pt] (6.521,6.961) -- (6.675,6.939);
\draw[draw=black!68, line width=0.17pt] (6.500,7.013) -- (6.296,7.042);
\draw[draw=black!65, line width=0.17pt] (6.524,7.001) -- (6.769,6.966);
\draw[draw=black!63, line width=0.17pt] (6.494,7.046) -- (6.373,7.063);
\draw[draw=black!60, line width=0.17pt] (6.538,7.039) -- (6.765,7.007);
\draw[draw=black!60, line width=0.17pt] (6.508,7.045) -- (6.253,7.081);
\draw[draw=black!55, line width=0.17pt] (6.535,7.052) -- (6.752,7.021);
\draw[draw=black!66, line width=0.17pt] (6.497,7.106) -- (6.280,7.137);
\draw[draw=black!64, line width=0.17pt] (6.541,7.095) -- (6.770,7.062);
\draw[draw=black!68, line width=0.17pt] (6.492,7.137) -- (6.324,7.161);
\draw[draw=black!57, line width=0.17pt] (6.545,7.142) -- (6.794,7.106);
\draw[draw=black!64, line width=0.17pt] (6.503,7.156) -- (6.339,7.180);
\draw[draw=black!67, line width=0.17pt] (6.562,7.153) -- (6.708,7.132);
\draw[draw=black!67, line width=0.17pt] (6.502,7.211) -- (6.269,7.245);
\draw[draw=black!55, line width=0.17pt] (6.562,7.209) -- (6.667,7.194);
\draw[draw=black!60, line width=0.17pt] (6.536,7.255) -- (6.325,7.285);
\draw[draw=black!63, line width=0.17pt] (6.559,7.262) -- (6.675,7.245);
\draw[draw=black!62, line width=0.17pt] (6.542,7.275) -- (6.374,7.299);
\draw[draw=black!68, line width=0.17pt] (6.576,7.268) -- (6.676,7.253);
\draw[draw=black!63, line width=0.17pt] (6.539,7.306) -- (6.416,7.261);
\draw[draw=black!61, line width=0.17pt] (6.570,7.321) -- (6.699,7.368);
\draw[draw=black!59, line width=0.17pt] (6.505,7.360) -- (6.288,7.281);
\draw[draw=black!65, line width=0.17pt] (6.539,7.364) -- (6.652,7.405);
\draw[draw=black!61, line width=0.17pt] (6.502,7.360) -- (6.308,7.289);
\draw[draw=black!52, line width=0.17pt] (6.535,7.386) -- (6.746,7.463);
\draw[draw=black!53, line width=0.17pt] (6.474,7.398) -- (6.279,7.327);
\draw[draw=black!55, line width=0.17pt] (6.514,7.412) -- (6.645,7.460);
\draw[draw=black!57, line width=0.17pt] (6.473,7.455) -- (6.284,7.386);
\draw[draw=black!52, line width=0.17pt] (6.526,7.464) -- (6.747,7.544);
\draw[draw=black!53, line width=0.17pt] (6.473,7.489) -- (6.297,7.425);
\draw[draw=black!65, line width=0.17pt] (6.516,7.478) -- (6.753,7.564);
\draw[draw=black!67, line width=0.17pt] (6.449,7.515) -- (6.205,7.427);
\draw[draw=black!52, line width=0.17pt] (6.489,7.521) -- (6.615,7.567);
\draw[draw=black!58, line width=0.17pt] (6.446,7.531) -- (6.262,7.464);
\draw[draw=black!66, line width=0.17pt] (6.477,7.539) -- (6.662,7.606);
\draw[draw=black!65, line width=0.17pt] (6.418,7.572) -- (6.181,7.486);
\draw[draw=black!65, line width=0.17pt] (6.472,7.591) -- (6.674,7.664);
\draw[draw=black!65, line width=0.17pt] (6.401,7.605) -- (6.229,7.543);
\draw[draw=black!58, line width=0.17pt] (6.444,7.638) -- (6.586,7.690);
\draw[draw=black!64, line width=0.17pt] (6.412,7.645) -- (6.271,7.594);
\draw[draw=black!58, line width=0.17pt] (6.447,7.656) -- (6.565,7.699);
\draw[draw=black!50, line width=0.17pt] (6.399,7.679) -- (6.185,7.601);
\draw[draw=black!64, line width=0.17pt] (6.434,7.685) -- (6.551,7.728);
\draw[draw=black!65, line width=0.17pt] (6.381,7.713) -- (6.253,7.666);
\draw[draw=black!66, line width=0.17pt] (6.417,7.718) -- (6.609,7.788);
\draw[draw=black!53, line width=0.17pt] (6.360,7.732) -- (6.183,7.667);
\draw[draw=black!56, line width=0.17pt] (6.411,7.755) -- (6.619,7.831);
\draw[draw=black!57, line width=0.17pt] (6.370,7.766) -- (6.164,7.691);
\draw[draw=black!57, line width=0.17pt] (6.399,7.769) -- (6.578,7.834);
\draw[draw=black!67, line width=0.17pt] (6.344,7.795) -- (6.239,7.757);
\draw[draw=black!50, line width=0.17pt] (6.370,7.815) -- (6.466,7.850);
\draw[draw=black!46, line width=0.17pt] (5.139,3.019) -- (4.985,3.066);
\draw[draw=black!47, line width=0.17pt] (5.183,3.007) -- (5.328,2.964);
\draw[draw=black!45, line width=0.17pt] (5.138,3.057) -- (4.986,3.103);
\draw[draw=black!43, line width=0.17pt] (5.186,3.044) -- (5.290,3.013);
\draw[draw=black!47, line width=0.17pt] (5.156,3.112) -- (5.069,3.137);
\draw[draw=black!49, line width=0.17pt] (5.194,3.127) -- (5.267,3.105);
\draw[draw=black!49, line width=0.17pt] (5.184,3.163) -- (5.070,3.197);
\draw[draw=black!43, line width=0.17pt] (5.201,3.147) -- (5.319,3.112);
\draw[draw=black!41, line width=0.17pt] (5.201,3.213) -- (5.067,3.253);
\draw[draw=black!38, line width=0.17pt] (5.244,3.196) -- (5.322,3.173);
\draw[draw=black!41, line width=0.17pt] (5.225,3.285) -- (5.090,3.325);
\draw[draw=black!52, line width=0.17pt] (5.248,3.263) -- (5.317,3.243);
\draw[draw=black!51, line width=0.17pt] (5.239,3.323) -- (5.099,3.364);
\draw[draw=black!46, line width=0.17pt] (5.264,3.297) -- (5.391,3.260);
\draw[draw=black!49, line width=0.17pt] (5.246,3.377) -- (5.163,3.402);
\draw[draw=black!42, line width=0.17pt] (5.278,3.363) -- (5.430,3.317);
\draw[draw=black!44, line width=0.17pt] (5.266,3.444) -- (5.193,3.466);
\draw[draw=black!46, line width=0.17pt] (5.306,3.428) -- (5.446,3.386);
\draw[draw=black!48, line width=0.17pt] (5.289,3.491) -- (5.140,3.536);
\draw[draw=black!38, line width=0.17pt] (5.312,3.487) -- (5.384,3.466);
\draw[draw=black!50, line width=0.17pt] (5.279,3.522) -- (5.126,3.568);
\draw[draw=black!44, line width=0.17pt] (5.339,3.529) -- (5.439,3.499);
\draw[draw=black!49, line width=0.17pt] (5.327,3.603) -- (5.240,3.629);
\draw[draw=black!52, line width=0.17pt] (5.364,3.599) -- (5.460,3.570);
\draw[draw=black!49, line width=0.17pt] (5.315,3.637) -- (5.163,3.682);
\draw[draw=black!41, line width=0.17pt] (5.374,3.653) -- (5.467,3.625);
\draw[draw=black!49, line width=0.17pt] (5.327,3.701) -- (5.223,3.733);
\draw[draw=black!40, line width=0.17pt] (5.370,3.678) -- (5.472,3.647);
\draw[draw=black!48, line width=0.17pt] (5.346,3.751) -- (5.271,3.774);
\draw[draw=black!44, line width=0.17pt] (5.392,3.725) -- (5.463,3.704);
\draw[draw=black!38, line width=0.17pt] (5.367,3.801) -- (5.244,3.838);
\draw[draw=black!49, line width=0.17pt] (5.426,3.801) -- (5.530,3.769);
\draw[draw=black!52, line width=0.17pt] (5.402,3.885) -- (5.332,3.906);
\draw[draw=black!44, line width=0.17pt] (5.423,3.856) -- (5.581,3.808);
\draw[draw=black!41, line width=0.17pt] (5.413,3.914) -- (5.305,3.946);
\draw[draw=black!42, line width=0.17pt] (5.433,3.927) -- (5.510,3.903);
\draw[draw=black!52, line width=0.17pt] (5.434,3.972) -- (5.344,3.999);
\draw[draw=black!40, line width=0.17pt] (5.458,3.975) -- (5.533,3.953);
\draw[draw=black!43, line width=0.17pt] (5.445,4.022) -- (5.359,4.026);
\draw[draw=black!45, line width=0.17pt] (5.482,4.027) -- (5.592,4.021);
\draw[draw=black!44, line width=0.17pt] (5.442,4.094) -- (5.306,4.102);
\draw[draw=black!44, line width=0.17pt] (5.466,4.072) -- (5.564,4.067);
\draw[draw=black!39, line width=0.17pt] (5.451,4.138) -- (5.295,4.147);
\draw[draw=black!46, line width=0.17pt] (5.465,4.121) -- (5.586,4.114);
\draw[draw=black!45, line width=0.17pt] (5.453,4.197) -- (5.329,4.204);
\draw[draw=black!48, line width=0.17pt] (5.476,4.196) -- (5.562,4.191);
\draw[draw=black!51, line width=0.17pt] (5.434,4.245) -- (5.313,4.252);
\draw[draw=black!40, line width=0.17pt] (5.479,4.238) -- (5.597,4.231);
\draw[draw=black!40, line width=0.17pt] (5.454,4.311) -- (5.382,4.315);
\draw[draw=black!41, line width=0.17pt] (5.492,4.310) -- (5.580,4.305);
\draw[draw=black!50, line width=0.17pt] (5.444,4.358) -- (5.336,4.364);
\draw[draw=black!52, line width=0.17pt] (5.489,4.331) -- (5.624,4.324);
\draw[draw=black!43, line width=0.17pt] (5.467,4.408) -- (5.301,4.418);
\draw[draw=black!44, line width=0.17pt] (5.497,4.405) -- (5.586,4.400);
\draw[draw=black!47, line width=0.17pt] (5.465,4.462) -- (5.375,4.467);
\draw[draw=black!52, line width=0.17pt] (5.492,4.462) -- (5.621,4.455);
\draw[draw=black!49, line width=0.17pt] (5.457,4.533) -- (5.310,4.541);
\draw[draw=black!51, line width=0.17pt] (5.497,4.544) -- (5.598,4.539);
\draw[draw=black!52, line width=0.17pt] (5.475,4.580) -- (5.344,4.588);
\draw[draw=black!50, line width=0.17pt] (5.512,4.574) -- (5.590,4.569);
\draw[draw=black!39, line width=0.17pt] (5.451,4.637) -- (5.362,4.642);
\draw[draw=black!48, line width=0.17pt] (5.503,4.654) -- (5.596,4.649);
\draw[draw=black!49, line width=0.17pt] (5.458,4.691) -- (5.320,4.699);
\draw[draw=black!46, line width=0.17pt] (5.518,4.706) -- (5.595,4.701);
\draw[draw=black!46, line width=0.17pt] (5.480,4.744) -- (5.391,4.749);
\draw[draw=black!52, line width=0.17pt] (5.516,4.739) -- (5.645,4.732);
\draw[draw=black!47, line width=0.17pt] (5.473,4.820) -- (5.335,4.828);
\draw[draw=black!49, line width=0.17pt] (5.528,4.814) -- (5.670,4.806);
\draw[draw=black!52, line width=0.17pt] (5.493,4.880) -- (5.326,4.889);
\draw[draw=black!52, line width=0.17pt] (5.523,4.866) -- (5.661,4.859);
\draw[draw=black!48, line width=0.17pt] (5.491,4.928) -- (5.375,4.957);
\draw[draw=black!47, line width=0.17pt] (5.511,4.924) -- (5.635,4.893);
\draw[draw=black!40, line width=0.17pt] (5.510,4.984) -- (5.404,5.011);
\draw[draw=black!41, line width=0.17pt] (5.534,4.993) -- (5.642,4.966);
\draw[draw=black!43, line width=0.17pt] (5.512,5.042) -- (5.381,5.075);
\draw[draw=black!50, line width=0.17pt] (5.553,5.051) -- (5.713,5.011);
\draw[draw=black!51, line width=0.17pt] (5.532,5.098) -- (5.371,5.138);
\draw[draw=black!45, line width=0.17pt] (5.580,5.073) -- (5.701,5.043);
\draw[draw=black!45, line width=0.17pt] (5.549,5.150) -- (5.465,5.171);
\draw[draw=black!44, line width=0.17pt] (5.584,5.117) -- (5.740,5.078);
\draw[draw=black!43, line width=0.17pt] (5.559,5.221) -- (5.437,5.252);
\draw[draw=black!41, line width=0.17pt] (5.598,5.194) -- (5.688,5.171);
\draw[draw=black!43, line width=0.17pt] (5.567,5.223) -- (5.406,5.263);
\draw[draw=black!43, line width=0.17pt] (5.612,5.221) -- (5.703,5.199);
\draw[draw=black!39, line width=0.17pt] (5.592,5.312) -- (5.482,5.339);
\draw[draw=black!39, line width=0.17pt] (5.616,5.293) -- (5.708,5.270);
\draw[draw=black!49, line width=0.17pt] (5.607,5.356) -- (5.469,5.390);
\draw[draw=black!50, line width=0.17pt] (5.621,5.371) -- (5.722,5.346);
\draw[draw=black!40, line width=0.17pt] (5.602,5.417) -- (5.509,5.440);
\draw[draw=black!42, line width=0.17pt] (5.647,5.418) -- (5.779,5.385);
\draw[draw=black!50, line width=0.17pt] (5.632,5.466) -- (5.516,5.495);
\draw[draw=black!41, line width=0.17pt] (5.666,5.469) -- (5.805,5.434);
\draw[draw=black!41, line width=0.17pt] (5.629,5.529) -- (5.560,5.546);
\draw[draw=black!38, line width=0.17pt] (5.684,5.510) -- (5.795,5.482);
\draw[draw=black!38, line width=0.17pt] (5.633,5.579) -- (5.513,5.609);
\draw[draw=black!43, line width=0.17pt] (5.675,5.562) -- (5.784,5.535);
\draw[draw=black!42, line width=0.17pt] (5.665,5.627) -- (5.524,5.663);
\draw[draw=black!44, line width=0.17pt] (5.716,5.626) -- (5.877,5.586);
\draw[draw=black!49, line width=0.17pt] (5.676,5.696) -- (5.589,5.718);
\draw[draw=black!45, line width=0.17pt] (5.716,5.673) -- (5.813,5.649);
\draw[draw=black!43, line width=0.17pt] (5.696,5.750) -- (5.552,5.786);
\draw[draw=black!47, line width=0.17pt] (5.726,5.754) -- (5.844,5.725);
\draw[draw=black!40, line width=0.17pt] (5.692,5.781) -- (5.613,5.800);
\draw[draw=black!45, line width=0.17pt] (5.746,5.779) -- (5.840,5.756);
\draw[draw=black!38, line width=0.17pt] (5.727,5.838) -- (5.582,5.874);
\draw[draw=black!45, line width=0.17pt] (5.743,5.830) -- (5.859,5.801);
\draw[draw=black!46, line width=0.17pt] (5.741,5.897) -- (5.622,5.927);
\draw[draw=black!48, line width=0.17pt] (5.771,5.890) -- (5.853,5.869);
\draw[draw=black!47, line width=0.17pt] (5.751,5.967) -- (5.643,5.994);
\draw[draw=black!47, line width=0.17pt] (5.782,5.945) -- (5.944,5.904);
\draw[draw=black!43, line width=0.17pt] (5.772,6.004) -- (5.669,6.030);
\draw[draw=black!49, line width=0.17pt] (5.787,6.016) -- (5.913,5.985);
\draw[draw=black!49, line width=0.17pt] (5.776,6.066) -- (5.689,6.088);
\draw[draw=black!45, line width=0.17pt] (5.794,6.052) -- (5.928,6.018);
\draw[draw=black!44, line width=0.17pt] (5.782,6.118) -- (5.648,6.152);
\draw[draw=black!38, line width=0.17pt] (5.823,6.097) -- (5.896,6.079);
\draw[draw=black!48, line width=0.17pt] (5.790,6.184) -- (5.645,6.220);
\draw[draw=black!48, line width=0.17pt] (5.835,6.157) -- (5.983,6.120);
\draw[draw=black!47, line width=0.17pt] (5.828,6.240) -- (5.739,6.262);
\draw[draw=black!52, line width=0.17pt] (5.860,6.228) -- (5.939,6.209);
\draw[draw=black!41, line width=0.17pt] (5.830,6.282) -- (5.715,6.311);
\draw[draw=black!45, line width=0.17pt] (5.869,6.276) -- (5.960,6.253);
\draw[draw=black!39, line width=0.17pt] (5.836,6.354) -- (5.717,6.384);
\draw[draw=black!41, line width=0.17pt] (5.866,6.338) -- (6.026,6.298);
\draw[draw=black!48, line width=0.17pt] (5.869,6.412) -- (5.708,6.452);
\draw[draw=black!51, line width=0.17pt] (5.895,6.410) -- (5.990,6.386);
\draw[draw=black!38, line width=0.17pt] (5.876,6.462) -- (5.755,6.492);
\draw[draw=black!39, line width=0.17pt] (5.900,6.453) -- (6.013,6.425);
\draw[draw=black!67, line width=0.17pt] (4.675,10.623) -- (4.815,10.754);
\draw[draw=black!60, line width=0.17pt] (4.639,10.596) -- (4.542,10.506);
\draw[draw=black!62, line width=0.17pt] (4.690,10.616) -- (4.788,10.708);
\draw[draw=black!50, line width=0.17pt] (4.680,10.561) -- (4.533,10.425);
\draw[draw=black!51, line width=0.17pt] (4.726,10.584) -- (4.836,10.686);
\draw[draw=black!59, line width=0.17pt] (4.701,10.556) -- (4.555,10.420);
\draw[draw=black!57, line width=0.17pt] (4.764,10.545) -- (4.851,10.626);
\draw[draw=black!56, line width=0.17pt] (4.739,10.499) -- (4.596,10.366);
\draw[draw=black!62, line width=0.17pt] (4.799,10.502) -- (4.900,10.595);
\draw[draw=black!64, line width=0.17pt] (4.780,10.471) -- (4.717,10.412);
\draw[draw=black!56, line width=0.17pt] (4.823,10.473) -- (4.903,10.548);
\draw[draw=black!61, line width=0.17pt] (4.797,10.446) -- (4.669,10.327);
\draw[draw=black!65, line width=0.17pt] (4.841,10.467) -- (4.984,10.600);
\draw[draw=black!58, line width=0.17pt] (4.811,10.428) -- (4.689,10.315);
\draw[draw=black!66, line width=0.17pt] (4.886,10.404) -- (5.002,10.512);
\draw[draw=black!61, line width=0.17pt] (4.857,10.404) -- (4.797,10.348);
\draw[draw=black!59, line width=0.17pt] (4.913,10.390) -- (5.041,10.508);
\draw[draw=black!60, line width=0.17pt] (4.896,10.350) -- (4.757,10.221);
\draw[draw=black!54, line width=0.17pt] (4.936,10.348) -- (5.083,10.484);
\draw[draw=black!55, line width=0.17pt] (4.907,10.343) -- (4.778,10.223);
\draw[draw=black!56, line width=0.17pt] (4.953,10.335) -- (5.040,10.415);
\draw[draw=black!58, line width=0.17pt] (4.943,10.294) -- (4.806,10.167);
\draw[draw=black!64, line width=0.17pt] (4.975,10.300) -- (5.112,10.426);
\draw[draw=black!63, line width=0.17pt] (4.961,10.275) -- (4.827,10.150);
\draw[draw=black!62, line width=0.17pt] (5.024,10.257) -- (5.136,10.360);
\draw[draw=black!57, line width=0.17pt] (4.975,10.235) -- (4.905,10.170);
\draw[draw=black!60, line width=0.17pt] (5.035,10.225) -- (5.108,10.292);
\draw[draw=black!66, line width=0.17pt] (5.025,10.204) -- (4.953,10.137);
\draw[draw=black!64, line width=0.17pt] (5.072,10.201) -- (5.209,10.338);
\draw[draw=black!59, line width=0.17pt] (5.057,10.159) -- (4.984,10.086);
\draw[draw=black!59, line width=0.17pt] (5.115,10.177) -- (5.183,10.246);
\draw[draw=black!58, line width=0.17pt] (5.081,10.148) -- (5.005,10.072);
\draw[draw=black!62, line width=0.17pt] (5.153,10.144) -- (5.276,10.267);
\draw[draw=black!58, line width=0.17pt] (5.128,10.108) -- (4.993,9.972);
\draw[draw=black!62, line width=0.17pt] (5.177,10.102) -- (5.240,10.165);
\draw[draw=black!63, line width=0.17pt] (5.122,10.072) -- (5.058,10.008);
\draw[draw=black!57, line width=0.17pt] (5.212,10.078) -- (5.288,10.154);
\draw[draw=black!63, line width=0.17pt] (5.180,10.063) -- (5.058,9.941);
\draw[draw=black!63, line width=0.17pt] (5.215,10.048) -- (5.311,10.144);
\draw[draw=black!66, line width=0.17pt] (5.184,10.015) -- (5.113,9.945);
\draw[draw=black!50, line width=0.17pt] (5.267,10.016) -- (5.360,10.109);
\draw[draw=black!55, line width=0.17pt] (5.241,10.000) -- (5.122,9.881);
\draw[draw=black!54, line width=0.17pt] (5.280,9.992) -- (5.391,10.103);
\draw[draw=black!50, line width=0.17pt] (5.244,9.979) -- (5.110,9.845);
\draw[draw=black!52, line width=0.17pt] (5.301,9.981) -- (5.365,10.045);
\draw[draw=black!54, line width=0.17pt] (5.284,9.947) -- (5.204,9.867);
\draw[draw=black!65, line width=0.17pt] (5.343,9.925) -- (5.444,10.026);
\draw[draw=black!54, line width=0.17pt] (5.314,9.893) -- (5.179,9.759);
\draw[draw=black!51, line width=0.17pt] (5.381,9.899) -- (5.451,9.968);
\draw[draw=black!66, line width=0.17pt] (5.338,9.875) -- (5.279,9.816);
\draw[draw=black!61, line width=0.17pt] (5.422,9.883) -- (5.554,10.014);
\draw[draw=black!56, line width=0.17pt] (5.390,9.834) -- (5.265,9.709);
\draw[draw=black!65, line width=0.17pt] (5.424,9.862) -- (5.564,10.002);
\draw[draw=black!60, line width=0.17pt] (5.392,9.813) -- (5.272,9.693);
\draw[draw=black!58, line width=0.17pt] (6.283,7.305) -- (6.264,7.417);
\draw[draw=black!50, line width=0.17pt] (6.279,7.254) -- (6.305,7.105);
\draw[draw=black!57, line width=0.17pt] (6.316,7.292) -- (6.291,7.438);
\draw[draw=black!67, line width=0.17pt] (6.321,7.255) -- (6.343,7.130);
\draw[draw=black!68, line width=0.17pt] (6.317,7.297) -- (6.288,7.462);
\draw[draw=black!60, line width=0.17pt] (6.344,7.256) -- (6.364,7.146);
\draw[draw=black!56, line width=0.17pt] (6.375,7.316) -- (6.351,7.447);
\draw[draw=black!66, line width=0.17pt] (6.370,7.277) -- (6.389,7.170);
\draw[draw=black!68, line width=0.17pt] (6.375,7.330) -- (6.354,7.450);
\draw[draw=black!52, line width=0.17pt] (6.401,7.284) -- (6.422,7.164);
\draw[draw=black!60, line width=0.17pt] (6.400,7.324) -- (6.384,7.412);
\draw[draw=black!67, line width=0.17pt] (6.427,7.293) -- (6.451,7.157);
\draw[draw=black!59, line width=0.17pt] (6.444,7.338) -- (6.430,7.418);
\draw[draw=black!63, line width=0.17pt] (6.444,7.305) -- (6.458,7.224);
\draw[draw=black!61, line width=0.17pt] (6.471,7.346) -- (6.452,7.451);
\draw[draw=black!67, line width=0.17pt] (6.484,7.306) -- (6.502,7.208);
\draw[draw=black!60, line width=0.17pt] (6.508,7.334) -- (6.479,7.497);
\draw[draw=black!61, line width=0.17pt] (6.504,7.299) -- (6.522,7.193);
\draw[draw=black!52, line width=0.17pt] (6.534,7.337) -- (6.513,7.454);
\draw[draw=black!64, line width=0.17pt] (6.522,7.298) -- (6.551,7.134);
\draw[draw=black!50, line width=0.17pt] (6.576,7.363) -- (6.551,7.506);
\draw[draw=black!61, line width=0.17pt] (6.585,7.302) -- (6.605,7.187);
\draw[draw=black!50, line width=0.17pt] (6.587,7.366) -- (6.561,7.512);
\draw[draw=black!60, line width=0.17pt] (6.600,7.311) -- (6.622,7.188);
\draw[draw=black!52, line width=0.17pt] (6.211,6.682) -- (6.234,6.793);
\draw[draw=black!46, line width=0.17pt] (6.193,6.627) -- (6.175,6.537);
\draw[draw=black!41, line width=0.17pt] (6.237,6.682) -- (6.254,6.767);
\draw[draw=black!51, line width=0.17pt] (6.210,6.639) -- (6.181,6.494);
\draw[draw=black!44, line width=0.17pt] (6.265,6.653) -- (6.291,6.783);
\draw[draw=black!48, line width=0.17pt] (6.255,6.616) -- (6.234,6.510);
\draw[draw=black!46, line width=0.17pt] (6.315,6.644) -- (6.336,6.751);
\draw[draw=black!43, line width=0.17pt] (6.307,6.606) -- (6.280,6.472);
\draw[draw=black!43, line width=0.17pt] (6.332,6.645) -- (6.352,6.743);
\draw[draw=black!52, line width=0.17pt] (6.318,6.604) -- (6.301,6.520);
\draw[draw=black!42, line width=0.17pt] (6.389,6.633) -- (6.407,6.724);
\draw[draw=black!47, line width=0.17pt] (6.367,6.601) -- (6.344,6.483);
\draw[draw=black!41, line width=0.17pt] (6.398,6.630) -- (6.418,6.727);
\draw[draw=black!40, line width=0.17pt] (6.406,6.611) -- (6.390,6.533);
\draw[draw=black!51, line width=0.17pt] (6.428,6.628) -- (6.454,6.755);
\draw[draw=black!43, line width=0.17pt] (6.445,6.585) -- (6.417,6.448);
\draw[draw=black!54, line width=0.17pt] (6.463,6.629) -- (6.492,6.775);
\draw[draw=black!50, line width=0.17pt] (6.454,6.576) -- (6.440,6.506);
\draw[draw=black!59, line width=0.17pt] (8.640,0.892) -- (8.792,0.896);
\draw[draw=black!66, line width=0.17pt] (8.644,0.941) -- (8.853,0.945);
\draw[draw=black!58, line width=0.17pt] (8.634,0.992) -- (8.784,0.995);
\draw[draw=black!62, line width=0.17pt] (8.639,1.007) -- (8.781,1.010);
\draw[draw=black!52, line width=0.17pt] (8.639,1.038) -- (8.842,1.042);
\draw[draw=black!52, line width=0.17pt] (8.633,1.097) -- (8.754,1.100);
\draw[draw=black!54, line width=0.17pt] (8.633,1.145) -- (8.750,1.148);
\draw[draw=black!57, line width=0.17pt] (8.623,1.177) -- (8.761,1.180);
\draw[draw=black!60, line width=0.17pt] (8.624,1.221) -- (8.755,1.224);
\draw[draw=black!56, line width=0.17pt] (8.644,1.249) -- (8.824,1.253);
\draw[draw=black!56, line width=0.17pt] (8.635,1.285) -- (8.808,1.289);
\draw[draw=black!54, line width=0.17pt] (8.638,1.311) -- (8.821,1.315);
\draw[draw=black!60, line width=0.17pt] (8.617,1.368) -- (8.811,1.372);
\draw[draw=black!51, line width=0.17pt] (8.639,1.421) -- (8.840,1.426);
\draw[draw=black!54, line width=0.17pt] (8.623,1.449) -- (8.724,1.452);
\draw[draw=black!56, line width=0.17pt] (8.640,1.477) -- (8.750,1.480);
\draw[draw=black!56, line width=0.17pt] (8.625,1.518) -- (8.764,1.521);
\draw[draw=black!66, line width=0.17pt] (8.633,1.561) -- (8.849,1.566);
\draw[draw=black!61, line width=0.17pt] (8.618,1.608) -- (8.814,1.613);
\draw[draw=black!53, line width=0.17pt] (8.618,1.623) -- (8.721,1.625);
\draw[draw=black!60, line width=0.17pt] (8.624,1.689) -- (8.798,1.693);
\draw[draw=black!59, line width=0.17pt] (8.636,1.719) -- (8.830,1.723);
\draw[draw=black!53, line width=0.17pt] (8.631,1.754) -- (8.775,1.758);
\draw[draw=black!64, line width=0.17pt] (8.630,1.783) -- (8.769,1.786);
\draw[draw=black!61, line width=0.17pt] (8.619,1.812) -- (8.838,1.800);
\draw[draw=black!64, line width=0.17pt] (8.625,1.867) -- (8.765,1.859);
\draw[draw=black!63, line width=0.17pt] (8.635,1.890) -- (8.835,1.879);
\draw[draw=black!51, line width=0.17pt] (8.620,1.960) -- (8.810,1.949);
\draw[draw=black!52, line width=0.17pt] (8.633,1.985) -- (8.737,1.979);
\draw[draw=black!65, line width=0.17pt] (8.626,2.019) -- (8.792,2.010);
\draw[draw=black!65, line width=0.17pt] (8.627,2.050) -- (8.800,2.041);
\draw[draw=black!59, line width=0.17pt] (8.643,2.101) -- (8.760,2.095);
\draw[draw=black!56, line width=0.17pt] (8.647,2.139) -- (8.818,2.129);
\draw[draw=black!60, line width=0.17pt] (8.627,2.162) -- (8.804,2.152);
\draw[draw=black!65, line width=0.17pt] (8.640,2.214) -- (8.792,2.205);
\draw[draw=black!57, line width=0.17pt] (8.632,2.254) -- (8.843,2.243);
\draw[draw=black!58, line width=0.17pt] (8.657,2.295) -- (8.850,2.284);
\draw[draw=black!60, line width=0.17pt] (8.642,2.316) -- (8.775,2.308);
\draw[draw=black!53, line width=0.17pt] (8.655,2.375) -- (8.835,2.365);
\draw[draw=black!51, line width=0.17pt] (8.663,2.403) -- (8.792,2.396);
\draw[draw=black!56, line width=0.17pt] (8.648,2.470) -- (8.760,2.464);
\draw[draw=black!54, line width=0.17pt] (8.657,2.491) -- (8.849,2.480);
\draw[draw=black!52, line width=0.17pt] (8.651,2.528) -- (8.755,2.522);
\draw[draw=black!60, line width=0.17pt] (8.667,2.576) -- (8.866,2.565);
\draw[draw=black!57, line width=0.17pt] (8.670,2.580) -- (8.849,2.571);
\draw[draw=black!57, line width=0.17pt] (8.660,2.622) -- (8.872,2.611);
\draw[draw=black!55, line width=0.17pt] (8.659,2.684) -- (8.779,2.678);
\draw[draw=black!57, line width=0.17pt] (8.665,2.710) -- (8.827,2.719);
\draw[draw=black!66, line width=0.17pt] (8.671,2.745) -- (8.824,2.753);
\draw[draw=black!59, line width=0.17pt] (8.657,2.819) -- (8.814,2.827);
\draw[draw=black!59, line width=0.17pt] (8.677,2.842) -- (8.789,2.849);
\draw[draw=black!51, line width=0.17pt] (8.674,2.877) -- (8.818,2.885);
\draw[draw=black!53, line width=0.17pt] (8.653,2.933) -- (8.823,2.943);
\draw[draw=black!51, line width=0.17pt] (8.652,2.965) -- (8.753,2.971);
\draw[draw=black!54, line width=0.17pt] (8.667,2.988) -- (8.805,2.995);
\draw[draw=black!58, line width=0.17pt] (8.654,3.017) -- (8.785,3.024);
\draw[draw=black!54, line width=0.17pt] (8.639,3.076) -- (8.840,3.087);
\draw[draw=black!58, line width=0.17pt] (8.660,3.094) -- (8.845,3.104);
\draw[draw=black!57, line width=0.17pt] (8.630,3.153) -- (8.838,3.164);
\draw[draw=black!66, line width=0.17pt] (8.657,3.187) -- (8.811,3.195);
\draw[draw=black!51, line width=0.17pt] (8.647,3.206) -- (8.753,3.211);
\draw[draw=black!51, line width=0.17pt] (8.642,3.253) -- (8.848,3.264);
\draw[draw=black!51, line width=0.17pt] (8.631,3.305) -- (8.806,3.315);
\draw[draw=black!50, line width=0.17pt] (8.641,3.351) -- (8.767,3.358);
\draw[draw=black!53, line width=0.17pt] (8.623,3.375) -- (8.786,3.384);
\draw[draw=black!56, line width=0.17pt] (8.624,3.423) -- (8.780,3.432);
\draw[draw=black!60, line width=0.17pt] (8.626,3.446) -- (8.796,3.456);
\draw[draw=black!64, line width=0.17pt] (8.631,3.500) -- (8.791,3.509);
\draw[draw=black!65, line width=0.17pt] (8.625,3.510) -- (8.759,3.517);
\draw[draw=black!53, line width=0.17pt] (8.629,3.570) -- (8.784,3.578);
\draw[draw=black!62, line width=0.17pt] (8.629,3.621) -- (8.790,3.603);
\draw[draw=black!65, line width=0.17pt] (8.618,3.636) -- (8.784,3.617);
\draw[draw=black!64, line width=0.17pt] (8.625,3.693) -- (8.835,3.669);
\draw[draw=black!63, line width=0.17pt] (8.624,3.702) -- (8.726,3.690);
\draw[draw=black!61, line width=0.17pt] (8.638,3.767) -- (8.783,3.751);
\draw[draw=black!63, line width=0.17pt] (8.639,3.794) -- (8.752,3.781);
\draw[draw=black!55, line width=0.17pt] (8.659,3.831) -- (8.802,3.815);
\draw[draw=black!57, line width=0.17pt] (8.645,3.877) -- (8.830,3.856);
\draw[draw=black!57, line width=0.17pt] (8.649,3.913) -- (8.835,3.892);
\draw[draw=black!58, line width=0.17pt] (8.671,3.957) -- (8.863,3.935);
\draw[draw=black!65, line width=0.17pt] (8.658,3.994) -- (8.806,3.977);
\draw[draw=black!63, line width=0.17pt] (8.656,4.027) -- (8.860,4.004);
\draw[draw=black!60, line width=0.17pt] (8.681,4.074) -- (8.856,4.054);
\draw[draw=black!65, line width=0.17pt] (8.677,4.111) -- (8.864,4.090);
\draw[draw=black!58, line width=0.17pt] (8.695,4.148) -- (8.869,4.128);
\draw[draw=black!63, line width=0.17pt] (8.675,4.177) -- (8.798,4.163);
\draw[draw=black!63, line width=0.17pt] (8.693,4.215) -- (8.902,4.191);
\draw[draw=black!50, line width=0.17pt] (8.683,4.240) -- (8.817,4.225);
\draw[draw=black!53, line width=0.17pt] (9.425,1.004) -- (9.299,1.023);
\draw[draw=black!51, line width=0.17pt] (9.461,1.008) -- (9.577,0.990);
\draw[draw=black!51, line width=0.17pt] (9.442,1.087) -- (9.330,1.104);
\draw[draw=black!43, line width=0.17pt] (9.484,1.058) -- (9.617,1.038);
\draw[draw=black!52, line width=0.17pt] (9.437,1.113) -- (9.272,1.137);
\draw[draw=black!51, line width=0.17pt] (9.497,1.106) -- (9.628,1.087);
\draw[draw=black!46, line width=0.17pt] (9.448,1.161) -- (9.290,1.185);
\draw[draw=black!47, line width=0.17pt] (9.497,1.181) -- (9.642,1.159);
\draw[draw=black!50, line width=0.17pt] (9.467,1.234) -- (9.358,1.250);
\draw[draw=black!40, line width=0.17pt] (9.508,1.234) -- (9.603,1.219);
\draw[draw=black!41, line width=0.17pt] (9.481,1.290) -- (9.397,1.303);
\draw[draw=black!44, line width=0.17pt] (9.499,1.274) -- (9.629,1.254);
\draw[draw=black!45, line width=0.17pt] (9.486,1.354) -- (9.322,1.379);
\draw[draw=black!44, line width=0.17pt] (9.531,1.368) -- (9.626,1.354);
\draw[draw=black!44, line width=0.17pt] (9.491,1.411) -- (9.355,1.432);
\draw[draw=black!53, line width=0.17pt] (9.526,1.420) -- (9.664,1.399);
\draw[draw=black!45, line width=0.17pt] (9.519,1.513) -- (9.370,1.536);
\draw[draw=black!45, line width=0.17pt] (9.538,1.502) -- (9.688,1.480);
\draw[draw=black!48, line width=0.17pt] (9.507,1.522) -- (9.332,1.548);
\draw[draw=black!47, line width=0.17pt] (9.542,1.528) -- (9.728,1.500);
\draw[draw=black!45, line width=0.17pt] (9.522,1.606) -- (9.443,1.618);
\draw[draw=black!42, line width=0.17pt] (9.567,1.603) -- (9.743,1.577);
\draw[draw=black!47, line width=0.17pt] (9.532,1.662) -- (9.414,1.679);
\draw[draw=black!40, line width=0.17pt] (9.557,1.677) -- (9.721,1.652);
\draw[draw=black!48, line width=0.17pt] (9.532,1.715) -- (9.398,1.735);
\draw[draw=black!41, line width=0.17pt] (9.586,1.710) -- (9.746,1.686);
\draw[draw=black!40, line width=0.17pt] (9.553,1.795) -- (9.418,1.815);
\draw[draw=black!40, line width=0.17pt] (9.584,1.795) -- (9.692,1.779);
\draw[draw=black!47, line width=0.17pt] (9.558,1.838) -- (9.454,1.853);
\draw[draw=black!50, line width=0.17pt] (9.587,1.858) -- (9.715,1.839);
\draw[draw=black!48, line width=0.17pt] (9.563,1.915) -- (9.424,1.935);
\draw[draw=black!54, line width=0.17pt] (9.606,1.903) -- (9.765,1.879);
\draw[draw=black!41, line width=0.17pt] (9.565,1.956) -- (9.415,1.979);
\draw[draw=black!40, line width=0.17pt] (9.623,1.950) -- (9.780,1.926);
\draw[draw=black!49, line width=0.17pt] (9.583,2.013) -- (9.467,2.002);
\draw[draw=black!46, line width=0.17pt] (9.606,2.019) -- (9.754,2.033);
\draw[draw=black!54, line width=0.17pt] (9.574,2.073) -- (9.478,2.065);
\draw[draw=black!42, line width=0.17pt] (9.607,2.078) -- (9.734,2.089);
\draw[draw=black!42, line width=0.17pt] (9.561,2.135) -- (9.455,2.126);
\draw[draw=black!52, line width=0.17pt] (9.603,2.117) -- (9.686,2.124);
\draw[draw=black!49, line width=0.17pt] (9.553,2.175) -- (9.420,2.163);
\draw[draw=black!44, line width=0.17pt] (9.616,2.170) -- (9.737,2.181);
\draw[draw=black!40, line width=0.17pt] (9.552,2.260) -- (9.446,2.250);
\draw[draw=black!41, line width=0.17pt] (9.592,2.266) -- (9.741,2.280);
\draw[draw=black!44, line width=0.17pt] (9.553,2.319) -- (9.369,2.302);
\draw[draw=black!42, line width=0.17pt] (9.578,2.313) -- (9.700,2.324);
\draw[draw=black!45, line width=0.17pt] (9.557,2.378) -- (9.461,2.369);
\draw[draw=black!47, line width=0.17pt] (9.575,2.368) -- (9.666,2.377);
\draw[draw=black!43, line width=0.17pt] (9.554,2.452) -- (9.386,2.436);
\draw[draw=black!54, line width=0.17pt] (9.592,2.459) -- (9.774,2.476);
\draw[draw=black!40, line width=0.17pt] (9.532,2.501) -- (9.384,2.487);
\draw[draw=black!53, line width=0.17pt] (9.580,2.510) -- (9.707,2.521);
\draw[draw=black!42, line width=0.17pt] (9.542,2.540) -- (9.414,2.528);
\draw[draw=black!42, line width=0.17pt] (9.576,2.554) -- (9.710,2.566);
\draw[draw=black!45, line width=0.17pt] (9.529,2.610) -- (9.392,2.598);
\draw[draw=black!49, line width=0.17pt] (9.567,2.612) -- (9.669,2.621);
\draw[draw=black!54, line width=0.17pt] (9.510,2.693) -- (9.416,2.685);
\draw[draw=black!42, line width=0.17pt] (9.569,2.684) -- (9.682,2.694);
\draw[draw=black!54, line width=0.17pt] (9.503,2.750) -- (9.381,2.739);
\draw[draw=black!51, line width=0.17pt] (9.538,2.742) -- (9.709,2.757);
\draw[draw=black!48, line width=0.17pt] (9.508,2.789) -- (9.362,2.776);
\draw[draw=black!48, line width=0.17pt] (9.552,2.799) -- (9.723,2.815);
\draw[draw=black!48, line width=0.17pt] (9.502,2.845) -- (9.417,2.837);
\draw[draw=black!43, line width=0.17pt] (9.553,2.836) -- (9.639,2.844);
\draw[draw=black!43, line width=0.17pt] (9.498,2.896) -- (9.352,2.883);
\draw[draw=black!50, line width=0.17pt] (9.532,2.918) -- (9.633,2.927);
\draw[draw=black!47, line width=0.17pt] (9.493,2.961) -- (9.329,2.946);
\draw[draw=black!47, line width=0.17pt] (9.542,2.966) -- (9.627,2.974);
\draw[draw=black!51, line width=0.17pt] (9.481,3.034) -- (9.314,3.018);
\draw[draw=black!47, line width=0.17pt] (9.519,3.047) -- (9.683,3.062);
\draw[draw=black!46, line width=0.17pt] (9.490,4.616) -- (9.344,4.624);
\draw[draw=black!41, line width=0.17pt] (9.507,4.605) -- (9.675,4.596);
\draw[draw=black!53, line width=0.17pt] (9.491,4.659) -- (9.359,4.667);
\draw[draw=black!48, line width=0.17pt] (9.514,4.674) -- (9.635,4.667);
\draw[draw=black!54, line width=0.17pt] (9.479,4.749) -- (9.376,4.755);
\draw[draw=black!45, line width=0.17pt] (9.514,4.748) -- (9.643,4.741);
\draw[draw=black!52, line width=0.17pt] (9.484,4.819) -- (9.317,4.829);
\draw[draw=black!53, line width=0.17pt] (9.535,4.818) -- (9.725,4.807);
\draw[draw=black!45, line width=0.17pt] (9.494,4.884) -- (9.329,4.893);
\draw[draw=black!41, line width=0.17pt] (9.539,4.866) -- (9.652,4.860);
\draw[draw=black!54, line width=0.17pt] (9.505,4.942) -- (9.338,4.952);
\draw[draw=black!45, line width=0.17pt] (9.549,4.941) -- (9.639,4.936);
\draw[draw=black!46, line width=0.17pt] (9.490,5.001) -- (9.406,5.006);
\draw[draw=black!51, line width=0.17pt] (9.544,5.023) -- (9.627,5.018);
\draw[draw=black!49, line width=0.17pt] (9.506,5.040) -- (9.347,5.049);
\draw[draw=black!51, line width=0.17pt] (9.553,5.051) -- (9.637,5.047);
\draw[draw=black!41, line width=0.17pt] (9.510,5.150) -- (9.427,5.155);
\draw[draw=black!54, line width=0.17pt] (9.561,5.159) -- (9.740,5.149);
\draw[draw=black!44, line width=0.17pt] (9.522,5.211) -- (9.355,5.221);
\draw[draw=black!54, line width=0.17pt] (9.554,5.199) -- (9.698,5.191);
\draw[draw=black!41, line width=0.17pt] (9.516,5.274) -- (9.419,5.279);
\draw[draw=black!51, line width=0.17pt] (9.544,5.268) -- (9.652,5.262);
\draw[draw=black!45, line width=0.17pt] (9.521,5.327) -- (9.379,5.335);
\draw[draw=black!47, line width=0.17pt] (9.557,5.322) -- (9.660,5.317);
\draw[draw=black!47, line width=0.17pt] (9.536,5.419) -- (9.348,5.429);
\draw[draw=black!50, line width=0.17pt] (9.553,5.422) -- (9.714,5.413);
\draw[draw=black!46, line width=0.17pt] (9.516,5.482) -- (9.327,5.492);
\draw[draw=black!40, line width=0.17pt] (9.580,5.500) -- (9.767,5.489);
\draw[draw=black!42, line width=0.17pt] (9.527,5.517) -- (9.423,5.495);
\draw[draw=black!48, line width=0.17pt] (9.559,5.542) -- (9.671,5.566);
\draw[draw=black!47, line width=0.17pt] (9.491,5.606) -- (9.327,5.571);
\draw[draw=black!52, line width=0.17pt] (9.538,5.627) -- (9.631,5.647);
\draw[draw=black!53, line width=0.17pt] (9.506,5.657) -- (9.407,5.636);
\draw[draw=black!46, line width=0.17pt] (9.544,5.685) -- (9.642,5.706);
\draw[draw=black!48, line width=0.17pt] (9.487,5.742) -- (9.337,5.710);
\draw[draw=black!50, line width=0.17pt] (9.505,5.755) -- (9.675,5.792);
\draw[draw=black!50, line width=0.17pt] (9.468,5.829) -- (9.288,5.790);
\draw[draw=black!44, line width=0.17pt] (9.490,5.823) -- (9.582,5.843);
\draw[draw=black!49, line width=0.17pt] (9.463,5.884) -- (9.313,5.852);
\draw[draw=black!52, line width=0.17pt] (9.495,5.897) -- (9.600,5.919);
\draw[draw=black!40, line width=0.17pt] (9.431,5.917) -- (9.254,5.879);
\draw[draw=black!50, line width=0.17pt] (9.483,5.934) -- (9.589,5.956);
\draw[draw=black!44, line width=0.17pt] (9.421,5.987) -- (9.317,5.965);
\draw[draw=black!53, line width=0.17pt] (9.462,6.003) -- (9.570,6.026);
\draw[draw=black!45, line width=0.17pt] (9.413,6.070) -- (9.237,6.032);
\draw[draw=black!44, line width=0.17pt] (9.430,6.086) -- (9.607,6.124);
\draw[draw=black!43, line width=0.17pt] (9.399,6.119) -- (9.219,6.081);
\draw[draw=black!44, line width=0.17pt] (9.422,6.120) -- (9.594,6.157);
\draw[draw=black!43, line width=0.17pt] (6.587,2.996) -- (6.507,3.030);
\draw[draw=black!43, line width=0.17pt] (6.615,3.076) -- (6.493,3.128);
\draw[draw=black!33, line width=0.17pt] (6.626,3.133) -- (6.510,3.183);
\draw[draw=black!44, line width=0.17pt] (6.678,3.220) -- (6.561,3.270);
\draw[draw=black!43, line width=0.17pt] (6.694,3.284) -- (6.632,3.311);
\draw[draw=black!36, line width=0.17pt] (6.721,3.342) -- (6.635,3.378);
\draw[draw=black!41, line width=0.17pt] (6.743,3.430) -- (6.638,3.475);
\draw[draw=black!41, line width=0.17pt] (6.797,3.522) -- (6.714,3.557);
\draw[draw=black!44, line width=0.17pt] (6.823,3.545) -- (6.765,3.571);
\draw[draw=black!34, line width=0.17pt] (6.857,3.655) -- (6.777,3.690);
\draw[draw=black!42, line width=0.17pt] (6.891,3.734) -- (6.827,3.755);
\draw[draw=black!37, line width=0.17pt] (6.896,3.796) -- (6.777,3.836);
\draw[draw=black!32, line width=0.17pt] (6.937,3.870) -- (6.865,3.893);
\draw[draw=black!41, line width=0.17pt] (6.957,3.921) -- (6.882,3.946);
\draw[draw=black!35, line width=0.17pt] (6.978,4.010) -- (6.909,4.033);
\draw[draw=black!33, line width=0.17pt] (6.999,4.067) -- (6.910,4.097);
\draw[draw=black!42, line width=0.17pt] (7.025,4.186) -- (6.896,4.229);
\draw[draw=black!39, line width=0.17pt] (7.068,4.213) -- (6.957,4.250);
\draw[draw=black!72, line width=0.19pt] (9.063,4.275) -- (9.032,4.428);
\draw[draw=black!64, line width=0.19pt] (9.055,4.215) -- (9.077,4.104);
\draw[draw=black!64, line width=0.19pt] (9.059,4.251) -- (9.029,4.399);
\draw[draw=black!64, line width=0.19pt] (9.085,4.215) -- (9.104,4.121);
\draw[draw=black!60, line width=0.19pt] (9.081,4.257) -- (9.066,4.336);
\draw[draw=black!67, line width=0.19pt] (9.108,4.224) -- (9.126,4.135);
\draw[draw=black!72, line width=0.19pt] (9.122,4.287) -- (9.095,4.422);
\draw[draw=black!60, line width=0.19pt] (9.150,4.252) -- (9.175,4.126);
\draw[draw=black!60, line width=0.19pt] (9.138,4.284) -- (9.120,4.374);
\draw[draw=black!64, line width=0.19pt] (9.153,4.243) -- (9.176,4.127);
\draw[draw=black!66, line width=0.19pt] (9.176,4.279) -- (9.153,4.396);
\draw[draw=black!56, line width=0.19pt] (9.188,4.240) -- (9.212,4.123);
\draw[draw=black!69, line width=0.19pt] (9.206,4.295) -- (9.184,4.404);
\draw[draw=black!62, line width=0.19pt] (9.213,4.250) -- (9.240,4.114);
\draw[draw=black!59, line width=0.19pt] (9.227,4.292) -- (9.197,4.444);
\draw[draw=black!67, line width=0.19pt] (9.245,4.261) -- (9.267,4.150);
\draw[draw=black!61, line width=0.19pt] (9.270,4.313) -- (9.246,4.434);
\draw[draw=black!69, line width=0.19pt] (9.277,4.252) -- (9.302,4.127);
\draw[draw=black!69, line width=0.19pt] (9.278,4.316) -- (9.251,4.452);
\draw[draw=black!67, line width=0.19pt] (9.284,4.263) -- (9.309,4.140);
\draw[draw=black!59, line width=0.19pt] (9.313,4.320) -- (9.282,4.473);
\draw[draw=black!61, line width=0.19pt] (9.315,4.274) -- (9.343,4.133);
\node[stiny, black!80, fill=white, inner sep=1pt] at (9.95,3.55) {Nebo:\\the hidden-ark\\tradition\\(2 Mach 2)};
\node[stiny, anchor=west, black!70, fill=white, inner sep=0.8pt] at (7.02,7.2) {Ebal};
\node[stiny, black!60, rotate=75, fill=white, fill opacity=0.75, text opacity=1, inner sep=0.8pt] at (5.05,2.8) {\emph{the hill country}};
% ---- the route of the ark ---------------------------------------------------
% (1) up from Sinai, east of the Salt Sea, to the plains of Moab
\draw[march] (9.7,-0.1)
  .. controls (9.5,0.9) and (9.25,2.2) .. (9.1,3.3)
  .. controls (9.0,4.0) and (8.9,4.4) .. (8.72,4.62);
\draw[-{Stealth[length=2mm]}, line width=0.68pt] (9.32,1.85) -- (9.26,2.15);
% across Jordan at Jericho
\draw[march] (8.72,4.62) .. controls (8.3,4.82) and (7.9,4.85) .. (7.56,4.88);
% (2) up the spine: Jericho to Sichem under Ebal, back down to Silo
\draw[march] (7.44,4.94) .. controls (7.1,5.6) and (6.8,6.3) .. (6.52,6.88);
\draw[march] (6.48,6.88) .. controls (6.52,6.6) and (6.54,6.35) .. (6.54,6.14);
% (3) out to battle: Silo west to Aphec
\draw[march] (6.44,6.08) .. controls (5.9,6.28) and (5.0,6.45) .. (4.5,6.4);
\draw[-{Stealth[length=2mm]}, line width=0.68pt] (5.3,6.42) -- (5.0,6.44);
% (4) captive among the five cities, then home by Bethsames
\draw[march] (4.32,6.28) .. controls (3.6,5.6) and (3.0,4.95) .. (2.78,4.32);
\draw[march] (2.82,4.14) .. controls (3.2,3.95) and (3.5,3.88) .. (3.8,3.9);
\draw[march] (3.9,4.0) -- (3.9,4.28);
\draw[march] (4.0,4.36) .. controls (4.25,4.32) and (4.45,4.26) .. (4.6,4.22);
\draw[march] (4.78,4.26) .. controls (5.0,4.38) and (5.15,4.46) .. (5.3,4.52);
% (6) the parked decades, then David's processions to Jerusalem
\draw[march] (5.5,4.54) .. controls (5.75,4.5) and (5.95,4.45) .. (6.08,4.42);
\draw[-{Stealth[length=2mm]}, line width=0.68pt] (5.95,4.46) -- (6.1,4.42);
% ---- sites ------------------------------------------------------------------
\fill (7.44,4.86) circle (0.065); \node[stiny, anchor=north, fill=white, inner sep=0.8pt] at (7.5,4.7) {Jericho};
\fill (6.48,6.96) circle (0.065); \node[stiny, anchor=east, fill=white, inner sep=0.8pt] at (6.36,7.0) {Sichem};
\fill (6.54,6.06) circle (0.065);
\draw[pline] (6.64,6.24) -- (6.78,6.46) -- (6.92,6.24);
\node[stiny, anchor=west, fill=white, inner sep=0.8pt] at (6.68,5.98) {Silo};
\fill (4.38,6.36) circle (0.065); \node[stiny, anchor=south, fill=white, inner sep=0.8pt] at (4.38,6.5) {Aphec};
\fill (2.7,4.2) circle (0.065);  \node[stiny, anchor=east, fill=white, inner sep=0.8pt] at (2.58,4.28) {Azotus};
\fill (3.9,3.9) circle (0.065);  \node[stiny, anchor=north east, fill=white, inner sep=0.8pt] at (3.85,3.8) {Geth};
\fill (3.9,4.38) circle (0.065); \node[stiny, anchor=south, fill=white, inner sep=0.8pt] at (3.72,4.52) {Accaron};
\fill (4.68,4.2) circle (0.065); \node[stiny, anchor=north, fill=white, inner sep=0.8pt] at (4.78,4.04) {Bethsames};
\fill (5.4,4.56) circle (0.065); \node[stiny, anchor=west, rotate=16, fill=white, inner sep=0.8pt] at (5.5,4.72) {Cariathiarim};
\fill (6.18,4.38) circle (0.075);
\draw[line width=0.4pt] (6.18,4.38) circle (0.13);
\node[stiny, anchor=north west, fill=white, inner sep=0.8pt] at (6.26,4.26) {\scshape Jerusalem};
% ---- stage numbers, in clear ground -----------------------------------------
\node[stage] at (9.65,1.7) {1};
\node[stage] at (6.9,5.5) {2};
\node[stage] at (5.2,6.85) {3};
\node[stage] at (2.6,3.55) {4};
\node[stage] at (4.55,3.3) {5};
\node[stage] at (6.0,3.8) {6};
% ---- from Sinai -------------------------------------------------------------
\node[stiny, anchor=east] at (9.05,0.2) {from Sinai (Num 10; 33)};
% ---- north arrow ------------------------------------------------------------
\draw[pline, -{Stealth[length=2.2mm]}] (10.55,10.6) -- (10.55,11.5);
\node[stiny] at (10.55,11.72) {N};
% ---- cartouche --------------------------------------------------------------
\draw[pline, fill=white] (0.4,10.05) rectangle (3.8,11.55);
\draw[pfine, draw=black!50] (0.5,10.15) rectangle (3.7,11.45);
\node[stiny, align=center] at (2.1,11.0) {\scshape the road of the ark\\[1pt]\upshape Sinai to Jerusalem};
\node[stiny, black!60] at (2.1,10.44) {a sketch, not a survey};
\end{scope}
\end{tikzpicture}
\sketchcaption{The road of the ark: (1) up from Sinai east of the Salt Sea
to the plains of Moab, in by Jordan and Jericho (Num 33; Josue 3--6);
(2) to Sichem under Ebal, then the long residence at Silo --- the tent marks
it; (3) out to battle and taken at Aphec (1 Kings 4); (4) seven months among
the Philistine cities, then home by Bethsames (1 Kings 5--6); (5) the parked
decades at Cariathiarim (1 Kings 7); (6) David's two processions to
Jerusalem (2 Kings 6). Sites stand at their traditional identifications
under the Douay's names, projected true to position but not to scale; where
an identification or an event's archaeology is disputed, the text says so at
the place. Nebo carries a tradition, not a finding.}
\end{sketchfig}


\subsection{The river and the stones}

Israel crosses at the worst possible moment, in the season of harvest when
the Jordan runs over his banks --- the narrator troubles to say so. The
people are commanded to follow the ark at an interval of some two thousand
cubits, something over half a mile, because they had never passed that way
before.\endnote{Josue 3:3--4, 15, Douay-Rheims (Challoner), read at
\url{https://drbo.org/chapter/06003.htm}; the flood-stage notice and the
two-thousand-cubit interval are paraphrased from the chapter as printed.} The
distance is already doing theology. A crowd keeps that kind of interval not
around what is useful but before what is sovereign: it is the space cleared
for a king, not the radius marked round a hazard. Then comes the title.
``Behold the ark of the covenant of the Lord of all the earth shall go
before you into the Jordan.''\endnote{Josue 3:11, Douay-Rheims (Challoner),
read at \url{https://drbo.org/chapter/06003.htm}, as quoted; repeated in
substance at 3:13, ``the ark of the Lord the God of the whole earth.''}
Nowhere else in the ark's whole story is it given this name, and it receives
it twice, both times at the river (3:11, 13). The God who is about to give
one small people one small land is announced at its border as owner of every
land. The conquest opens with a claim that relativizes it. The crossing
itself is told as liturgy rather than stratagem: the priests bear the ark
down into the flood, the waters stand and are heaped upstream, ``swelling up
like a mountain,'' while the bearers ``stood girded upon the dry ground in
the midst of the Jordan'' and the nation passed over.\endnote{Josue 3:13,
15--17, Douay-Rheims (Challoner), read at
\url{https://drbo.org/chapter/06003.htm}; quoted phrases as printed there:
the waters ``shall stand together upon a heap'' (3:13), ``swelling up like a
mountain,'' ``stood girded upon the dry ground in the midst of the Jordan''
(3:15--17).}

What Israel is given to keep from the marvel is not a relic but a catechism.
Twelve stones are taken from the riverbed where the priests' feet had stood
and set up at Gilgal. Josue sets up twelve more in the midst of the Jordan
itself, and the narrator adds that they were there still in his own day. The
stones exist to provoke a question from children and to oblige an answer
from fathers: when your sons ask what these stones mean, you shall tell them
what the Lord did at this river. When the priests come up out of the
channel, the waters return to their place.\endnote{Josue 4:6--9, 18, 21--24,
paraphrased at the locus (wording checked in a modern edition, not
re-anchored in Douay for quotation). Gilgal itself has never been
identified; candidate sites northeast of Jericho remain unconfirmed, and
this study says only ``near Jericho, site unknown.''} The first thing the
ark does in the land, after the wonder, is to leave behind a discipline of
memory. The wonder serves the telling, not the telling the wonder.

\subsection{The wall and the mountain of the curse}

At Jericho the text is exact about the order of march, and the order is the
argument. Armed men go first; then seven priests with seven trumpets ---
trumpets the Douay names by their festal office, those ``used in the
jubilee''; then the ark; and ``the rest of the common people followed the
ark.'' So arranged, ``the ark of the Lord went about the city once a day''
for six days, and seven times on the seventh. The narrative gives the wall's
fall no mechanism beyond the circuits, the trumpets, and the
shout.\endnote{Josue 6:4, 6--13, Douay-Rheims (Challoner), read at
\url{https://drbo.org/chapter/06006.htm}; quoted phrases as printed: the
trumpets ``used in the jubilee'' (in the opening instructions, 6:4--8), ``the
ark of the Lord went about the city once a day'' (6:11), ``the rest of the
common people followed the ark'' (6:13).} The army is present, but as
escort. The procession is shaped from its center outward, and the center is
the ark. Jericho is not besieged. It is circled, as if in pilgrimage, until
it is not there. The walls fell to a liturgy.

The mound itself is not in doubt: Jericho is Tell es-Sultan, fixed by its
name, its spring, and a century of excavation. The archaeology of the fall
is disputed --- the fortified Bronze Age city was found destroyed well
before any plausible date for Josue --- and the dispute must simply be
reported.\endnote{Site identification and the state of the
question are reported from the standard excavation literature (Sellin and
Watzinger, Garstang, Kenyon): the identification of Tell es-Sultan is
secure; Kathleen Kenyon's finding that the fortified Bronze Age city was
destroyed long before any plausible conquest horizon is the crux of the
event-archaeology dispute. Reported at one remove; not re-examined for this
study.} And when Israel is beaten at Ai, the ark shows its other face: Josue
rends his garments and lies on his face before the ark of the Lord until
evening.\endnote{Josue 7:6, paraphrased at the locus.} The object at whose
passage walls fell is also the place where defeat is carried and asked
about. In both fortunes, the ark is where Israel goes to be before God.

Then, with the conquest scarcely begun, the campaign stops for a covenant.
At Shechem, between the two mountains, all Israel stands in two halves ---
the stranger as well as the homeborn. One half faces Gerizim for the
blessing, the other faces Ebal for the curse. Between them, facing both, the
Levitical priests carry the ark, and the words of the law, blessing and
curse together, are read over the whole assembly: the liturgy Moses had
commanded in Deuteronomy 27, performed around the box that holds the
tables.\endnote{Josue 8:30--35, paraphrased at the locus (8:33 checked in a
modern edition this study's research records identify; not re-anchored in
Douay for quotation); the commanded rite is Deuteronomy 27. The site of
Shechem is securely Tell Balata; the structure on Mount Ebal excavated by
Adam Zertal and read by him as Josue's altar remains sharply contested, and
the associated ``curse tablet'' claim of 2019--2022 has not won peer
acceptance --- both reported, neither relied on.} The scene is the plainest
exposition of the ark in the whole conquest. Jericho showed what the ark's
presence can do; Ebal shows what the ark is for. The nation arranged around
it is not an army but a congregation, and what issues from the center is not
power but the reading of the law. Covenant, not conquest, is the content of
the box.

\subsection{Silo}

When the land is divided, the tent of meeting is pitched at Silo --- Shiloh,
in the familiar spelling --- in the hill country of Ephraim (Josue 18:1).
There the ark stops for the longest residence of its story, from the days of
Josue to the childhood of Samuel. Through the whole age of the Judges it is
glimpsed exactly once: at Bethel, during the civil war against Benjamin,
with Phinees the son of Eleazar ministering before it. That chronological
notice is commonly read as setting the book's grim epilogue early in the
period.\endnote{Judges 20:27--28, paraphrased at the locus (the Douay's
temporal marker there is ``At that time''); the inference
from the chronological notice to an early setting for the book's closing
episodes is an interpretive commonplace, reported as such. The pitching of
the tent at Silo is Josue 18:1; the yearly feast at Silo is Judges 21:19.}
Otherwise, silence; and here silence means residence. By the time the story
resumes, Silo is a settled sanctuary: a yearly feast, a hereditary
priesthood, pilgrim families going up to sacrifice, a lamp burning through
the night --- and a narrative that has quietly begun to call the tent a
temple.

\begin{witnessbox}{1 Kings 3:3 (= 1 Samuel 3:3), Douay-Rheims}
Before the lamp of God went out, Samuel slept in the temple of the Lord,
where the ark of God was.
\end{witnessbox}

The verse earns its frame by what happens next in the chapter: in that room,
in the dark, the Lord calls a child by name, and the child learns to answer.
This is the promise of Exodus 25:22 being kept. The God who said he would
speak from above the propitiatory does speak, in the house where the ark is,
to a boy lying near it.\endnote{1 Kings 3:3 (= 1 Samuel 3:3), Douay-Rheims
(Challoner), read at \url{https://drbo.org/chapter/09003.htm}, as quoted; the
call follows in 3:4--14. The meeting-promise is Exodus 25:22, treated in the
Sinai section above.} After a river stopped and a wall thrown down, the
ark's deepest work in the land is the quietest: it marks the place where God
speaks. That is what it was built for, and at Silo, for once, it is allowed
simply to do it.

What became of Silo the Scripture never narrates. It remembers it, twice,
from a distance of centuries. The psalmist, recounting Israel's
rebellions, says that God ``put away the tabernacle of
Silo''\endnote{Psalm 77:60, Douay numbering (= Psalm 78:60), Douay-Rheims
(Challoner), read at \url{https://drbo.org/chapter/21077.htm}; the fragment
is quoted as printed there. Verses 60--61 are block-quoted, with their
textual disclosure, where the pattern of the journey is drawn below.} --- a
verse that must be heard in full when the time comes to weigh what
departed from Israel with the ark.
And Jeremias, preaching against a temple grown complacent: ``Go ye to my
place in Silo, where my name dwelt from the beginning: and see what I did to
it for the wickedness of my people Israel.''\endnote{Jeremias 7:12,
Douay-Rheims (Challoner), read at \url{https://drbo.org/chapter/28007.htm},
as quoted; the threat is turned on the temple itself at 7:14, and the sermon
is reprised in Jeremias 26.} And here, for once in this story, the dirt has
something to lay beside the text. Shiloh is Khirbet Seilun, an
identification as secure as any in the archaeology of Israel. The excavated
early Iron Age settlement there ends in a heavy burn layer: ash, collapsed
brick, scorched store-jars, and a great deposit of burned animal bones that
bespeaks a place of sacrifice. The layer is conventionally dated near the
middle of the eleventh century before Christ --- the right horizon for the
catastrophe the next section will reach.\endnote{The identification rests on
the survival of the name, the precise directions of Judges 21:19, and a
century of excavation. Excavation sequence: the Danish expedition under Hans
Kjaer, 1926--1932 (with a season in 1963); Israel Finkelstein for Tel Aviv
University, 1981--1984; renewed soundings from 1988; the Associates for
Biblical Research under Scott Stripling since 2017. Sequence per the site
reports read for this study,
\url{https://emekshaveh.org/en/wp-content/uploads/2013/07/13-Tel-Shiloh-Eng-03.pdf}
and
\url{https://biblearchaeology.org/50-the-shiloh-excavations/4345-research-on-shiloh};
the burn-layer particulars, the mid-eleventh-century date, and the cultic
reading of the bone deposit are the excavators' published account, reported
at one remove.} The convergence should be stated exactly, because it is one
of the few places in this whole history where text and stratum genuinely
meet. No inscription names the destroyer. The attribution to the
Philistines is an inference from the narrative, not a find, and the date is
conventional. What may honestly be said is this: the one sanctuary
Scripture remembers as destroyed is the one whose soil shows a destruction,
at about the remembered time. The two witnesses do not prove each other.
They face each other, and neither blinks.

But burn layer and psalm alike lie far downstream of the point the story has
reached. In narrative time Silo is at peace, and the peace is real: the lamp
burns, the feast comes round each year, a family climbs the hill country with
its offering, a child sleeps beside the ark and is learning to answer when he
is called. Only, the same chapters that show us the sleeping child show us
the waking priests: the sons of Heli, who treat the Lord's sacrifice as their
own larder, and against whose father's house a man of God has already
pronounced sentence (1 Kings 2).\endnote{The conduct of the sons of Heli and
the oracle of the man of God against the house of Heli: 1 Kings 2 (= 1 Samuel
2), paraphrased at the chapter.} The reader is not asked to guess; the text
has told him. The ark has rested at Silo through the longest age of its
history. The next time it goes out, Ophni and Phinees will go with it; and it
will not come back.
